% =====================================================================
%  Conjectures on finding a k-way collision with bounded memory.
%  Standalone; compiles with pdflatex, no external files needed.
%
%  LABELLING.  Every theorem, conjecture and problem below is named by a
%  code, and the code is also its \label, so \cref{RS4} prints
%  "Conjecture RS4".  A code is <model><mode><rank>:
%
%      model   U = unconstrained (no space bound)
%              R = register machine (state holds points only)
%              B = branching program (state arbitrary, width bounded)
%              G = general adversary (state arbitrary, cumulative
%                  memory charged)
%      mode    S = sequential (one query per step)
%              P = parallel (q queries per round)
%      rank    index within the class, increasing in logical strength
%      -k      suffix marking the general-k counterpart of a k=3 code
%
%  CONCORDANCE with the previous draft, in which the conjectures were
%  C1-C7 and G1-G5 and the ladder rungs were R1-R6, B1-B5, M1-M4,
%  P1-P3, Q1-Q3:
%
%      C1 -> RS2     C2 -> RS3     C3 -> RS4     C4 -> BS4    C5 -> GS3
%      C6 -> GP3     C7 -> BP3
%      G1 -> RS2-k   G2 -> RS3-k   G4 -> BS4-k   G5 -> GS3-k
%      Problem (register curve, k >= 4)  ->  Problem RS4-k
%      R1 -> RS0 (now a theorem with a proof)    R2 -> RS1   R6 -> RS5
%      B1 -> BS0   B2 -> BS1   B3 -> BS2   B4 -> BS3   B5 -> BS4
%      M1 -> GS1   M2 -> GS2   M3 -> GS3   M4 -> GS4
%      P1 -> GP1   P2 -> GP2   P3 -> GP3
%      Q1 -> BP1   Q2 -> BP2   Q3 -> BP3
%      new formal statements: GP3-k and BP3-k, which the previous draft
%      gave only in prose, and RS5, BS0-BS3, GS1-GS2, GS4, GP1-GP2,
%      BP1-BP2, which it gave only as ladder rungs.
%
%      SECOND REVISION.  Five codes are now theorems and one Problem has
%      become a Conjecture, after Dinur's streaming switching lemma
%      (eprint 2019/413), Beame-Kornerup (ICALP 2023), Guo-Nan-Yao
%      (Asiacrypt 2025) and Hao-Huang-Liu (arXiv 2602.23763):
%          BS0  now a theorem, and strengthened to the whole curve
%          RS5' new: RS5 proved for lifting machines
%          BP0  new: BS0 for oblivious parallel programs
%          GP0  new: the same in cumulative memory
%          BS4-C, GS3-C new: the C-fold statements, where the known
%                       technique applies
%          RS4-k  Problem -> Conjecture; the upper bound is now attained
%
%      THIRD REVISION.  Motivation added, the register model grounded in
%      Maurer's abstract models, and three neighbouring models that do
%      not exhibit the phenomenon recorded in a new remark (advice,
%      decision, comparison-based).  No code changed meaning.
%
%  The ladders are no longer a separate section: each class now presents
%  its own ladder in rank order, Section 5 indexes all codes, and
%  Section 6 orders them by how hard they are to resolve, and Appendix A
%  gives all of the models as code-based games.
%
%      FOURTH REVISION.  Appendix A added: the classes as restrictions on
%      the type of one transition function, the execution as a game
%      against a lazily sampled oracle, the two transforms that
%      segmentation proofs use, and four conventions that the prose model
%      boxes leave implicit.  The feedback channel of the register model
%      is now specified in Section 2.1 rather than left to the phrase
%      "what it was last told"; see Remark A on why the choice matters.
% =====================================================================
\documentclass[11pt]{article}
\usepackage[margin=1.15in]{geometry}
\usepackage{amsmath,amssymb,amsthm}
\usepackage{booktabs}
\usepackage{enumitem}
\usepackage[most]{tcolorbox}
\usepackage[colorlinks,linkcolor=blue!55!black,citecolor=blue!60!black,urlcolor=blue!60!black,bookmarksnumbered,bookmarksopen]{hyperref}
\usepackage[capitalise,nameinlink,nosort]{cleveref}

% --- theorem environments -------------------------------------------
% Coded statements (about the collision problem) print their code as
% their number; uncoded ones (about the models) are numbered normally.
\theoremstyle{plain}
\newtheorem{innerthm}{Theorem}
\newtheorem{innerconj}{Conjecture}
\newtheorem{proposition}{Proposition}
\newtheorem{corollary}{Corollary}
\newtheorem{lemma}{Lemma}
\theoremstyle{definition}
\newtheorem{innerprob}{Problem}
\newtheorem{remark}{Remark}
\newtheorem{definition}{Definition}
\newtheorem{convention}{Convention}

\newenvironment{cthm}[1]%
  {\renewcommand{\theinnerthm}{#1}%
   \def\theHinnerthm{cthm.\arabic{innerthm}}\innerthm}{\endinnerthm}
\newenvironment{conj}[1]%
  {\renewcommand{\theinnerconj}{#1}%
   \def\theHinnerconj{conj.\arabic{innerconj}}\innerconj}{\endinnerconj}
\newenvironment{cprob}[1]%
  {\renewcommand{\theinnerprob}{#1}%
   \def\theHinnerprob{cprob.\arabic{innerprob}}\innerprob}{\endinnerprob}

\crefname{innerthm}{Theorem}{Theorems}
\Crefname{innerthm}{Theorem}{Theorems}
\crefname{innerconj}{Conjecture}{Conjectures}
\Crefname{innerconj}{Conjecture}{Conjectures}
\crefname{innerprob}{Problem}{Problems}
\Crefname{innerprob}{Problem}{Problems}
\crefname{proposition}{Proposition}{Propositions}
\Crefname{proposition}{Proposition}{Propositions}
\crefname{corollary}{Corollary}{Corollaries}
\Crefname{corollary}{Corollary}{Corollaries}
\crefname{remark}{Remark}{Remarks}
\Crefname{remark}{Remark}{Remarks}
\crefname{lemma}{Lemma}{Lemmas}
\Crefname{lemma}{Lemma}{Lemmas}
\crefname{definition}{Definition}{Definitions}
\Crefname{definition}{Definition}{Definitions}
\crefname{convention}{Convention}{Conventions}
\Crefname{convention}{Convention}{Conventions}
\crefname{figure}{Figure}{Figures}
\Crefname{figure}{Figure}{Figures}

\newcommand{\NN}{[N]}
\newcommand{\advers}{\mathcal{A}}
\newcommand{\cmcs}{\mathrm{cmc}^{*}}
\newcommand{\eCMCs}{\mathrm{eCMC}^{*}}
\newcommand{\CMCs}{\mathrm{CMC}^{*}}
\newcommand{\Advance}{\mathsf{Advance}}
\newcommand{\Next}{\mathsf{Next}}
\newcommand{\Copy}{\mathsf{Copy}}
\newcommand{\Update}{\mathsf{Update}}
\newcommand{\cls}[1]{\textsf{#1}}   % a class, e.g. \cls{RS}
\newcommand{\Halt}{\mathsf{Halt}}
\newcommand{\getsr}{\mathrel{\xleftarrow{\raisebox{-0.15ex}{\tiny\$}}}}
\newcommand{\qry}{\mathsf{qry}}
\newcommand{\nxt}{\mathsf{nxt}}
\newcommand{\outp}{\mathsf{out}}
\newcommand{\stp}{\mathsf{step}}
\newcommand{\Exec}{\mathrm{EXEC}}
\newcommand{\Col}{\mathrm{3COL}}
\newcommand{\RO}{\mathsf{RO}}
\newcommand{\Fin}{\mathsf{Fin}}
\newcommand{\Ver}{\mathsf{Ver}}
\newcommand{\Adv}{\mathbf{Adv}}

% pseudocode lists and game boxes, used in Appendix A
\newlist{pc}{enumerate}{1}
\setlist[pc]{label=\scriptsize\arabic*:,ref=\arabic*,leftmargin=2.0em,
  itemsep=1.5pt,topsep=3pt,parsep=0pt,partopsep=0pt}
\newcommand{\gname}[1]{\underline{\smash{#1}}\vspace{2pt}}

\newtcolorbox{cbox}{colback=black!3,colframe=black!55,boxrule=0.5pt,
  arc=2pt,left=7pt,right=7pt,top=6pt,bottom=6pt}
\newtcolorbox{gbox}{colback=black!2,colframe=black!60,boxrule=0.5pt,
  arc=2pt,left=6pt,right=6pt,top=5pt,bottom=5pt}

\title{Conjectures on Finding a $k$-Way Collision\\ with Bounded Memory}
\author{}
\date{}

\begin{document}
\maketitle
\vspace{-4em}

\begin{abstract}
\noindent
A random function $f\colon \NN \to \NN$ has three-way collisions in
abundance, and $\Theta(N^{2/3})$ oracle queries find one.  Every algorithm
achieving that count stores a table of size $N^{1/3}$, and the best known
algorithm using no table needs $\Theta(N)$ queries.  Whether the table is
necessary is open.  This note states the conjectures to that effect as a
single labelled family.  A code names each one: the letters give its
\emph{class}, a memory model (register machine, branching program, or
general adversary charged for cumulative memory) crossed with a mode
(sequential or parallel); the index gives its \emph{rank} within the
class, ordered by logical strength; and a suffix \textsf{-k} marks the
general-$k$ counterpart.  Every code names a formal statement.  Five are
theorems, proved here or assembled from the literature, and
\cref{sec:route} orders the rest by how hard they look to resolve, which is
not the order of logical strength.  All of the conjectures are \emph{false}
for two-way collisions, which is the one thing every proof attempt must
respect.
\end{abstract}

% ---------------------------------------------------------------------
\section{Setup, and what is unconditional}
% ---------------------------------------------------------------------

Throughout, $N \geq 2$ and $f \colon \NN \to \NN$ is a uniformly random
function, available only through an oracle answering $x \mapsto f(x)$.  A
\emph{$k$-collision} is a set of $k$ distinct points of $\NN$ with the same
image under $f$; a \emph{$3$-collision} is the case $k=3$.  Asymptotic
notation $\tilde O, \tilde\Omega, \tilde\Theta$ suppresses factors
polylogarithmic in $N$.  Three facts are theorems and frame everything
below.

\begin{enumerate}[label=(F\arabic*),ref=F\arabic*,leftmargin=*]\itemsep2pt
\item\label{F1} \emph{The query threshold is $N^{2/3}$.}  Among $q$
  queried points the expected number of $3$-collisions is
  $\binom{q}{3}N^{-2}$, so $q = \Theta(N^{2/3})$ is necessary and
  sufficient for constant success probability, with the sharp constant
  $(3!)^{1/3}$ \cite{STKT06}.  Quantitatively, an algorithm making $T$
  queries and then naming three points succeeds with probability
  $O\bigl((T^{3}+1)N^{-2}\bigr)$, by Markov's inequality together with a
  term for naming points it never queried.  This holds for \emph{any}
  algorithm and no conjecture below contradicts it.
\item\label{F2} \emph{Without a table, $O(N)$ suffices.}  Find a
  $2$-collision $\{u,v\}$ by cycle-finding, retain $u$, $v$ and
  $y=f(u)$, then sample fresh points until one maps to $y$.  This uses
  $O(\log N)$ bits of memory and $\Theta(N)$ queries.
\item\label{F3} \emph{At $k=2$ memory is free.}  Pollard's rho
  method~\cite{Pollard75}, with Floyd's or Brent's cycle
  detection~\cite{Brent80}, finds a $2$-collision in $\Theta(\sqrt N)$
  queries with $O(\log N)$ bits, because a random trajectory closes on
  itself after $\Theta(\sqrt N)$ steps~\cite{FO90} and the cycle-entry
  vertex then has two discovered preimages.  That count is optimal even
  with unbounded memory.  Adaptivity is what buys it: on a query sequence
  fixed in advance, no algorithm with $O(\log N)$ bits comes near
  $\sqrt N$, a separation quantified by the streaming bounds
  of~\cite{JT19,Din20b} and pursued for preprocessing adversaries
  in~\cite{DKM25}.
\end{enumerate}

Between~\ref{F1} and~\ref{F2} lies a gap of $N^{1/3}$, and the conjectures
assert that memory is what closes it.  Interpolating, Joux and
Lucks~\cite{JL09} achieve $T = \tilde O(N/S)$ with a table of $S$
collisions for every $S \leq N^{1/3}$, so every conjecture below is
tight if true.  Fact~\ref{F3} is the constraint: \emph{every} conjecture
here is false at $k=2$, so no proof of any of them can be indifferent to
the collision order.

Three reasons to want the gap explained.  Memory-hard proof of work rests
on claims of this shape: Equihash~\cite{BK16} fixes its parameters from the
presumed time--memory behaviour of Wagner's generalised birthday
algorithm~\cite{Wag02}, which is a $k$-collision search, and the cumulative
measure of \cref{sec:cmc} was introduced to make such claims
precise~\cite{AS15}.  Second, whether $k$-collision resistance is strictly
weaker than collision resistance has different answers in different models:
in the standard model it is barely weaker, since collision resistance
follows from $k$-collision resistance for every constant
$k$~\cite{RV22,BT24}, and against Merkle--Damg\aa{}rd hashing with $S$ bits
of preprocessed advice the $m$-way advantage falls below the $2$-way one by
a factor $m$ and no more~\cite{Aks24}.  The conjectures here say that
against bounded \emph{working} memory the separation is instead a power of
$N$.  Third, $k=3$ at $S=O(1)$ is the smallest instance for which no proved
bound mentions memory at all.

\subsection{The label alphabet}\label{sec:labels}

A statement about this problem is fixed by three choices: what the
algorithm's memory may contain, whether it queries one point at a time or
$q$ at once, and how much it is asked to prove.  The first two choices name
a \emph{class} and are recorded in the letters of a code; the third is
recorded in the index.

\begin{center}
\small
\begin{tabular}{@{}ll@{}}
\toprule
\multicolumn{2}{@{}l}{\emph{first letter: what the memory may contain}}\\
\cls{U} & unconstrained: no bound on memory \\
\cls{R} & register machine: the state holds points, and nothing else \\
\cls{B} & branching program: the state is arbitrary, the width bounded \\
\cls{G} & general adversary: the state is arbitrary, cumulative memory charged \\
\addlinespace[3pt]
\multicolumn{2}{@{}l}{\emph{second letter: how queries are issued}}\\
\cls{S} & sequential: one query per step \\
\cls{P} & parallel: $q$ queries per round, $r$ rounds, $T = rq$ \\
\bottomrule
\end{tabular}
\end{center}

\medskip\noindent
So \cls{RS} is the sequential register machine of
\cref{sec:reg}, \cls{BP} the parallel branching program of
\cref{sec:GP}, and \cls{GP} the general parallel adversary measured
by cumulative memory.  Within a class the index increases with logical
strength, so \textsf{RS1} is the weakest statement of class \cls{RS} and
\textsf{RS5} the strongest; \cref{p:ladders} proves the ordering.  The
suffix \textsf{-k} marks the general-$k$ counterpart of a $k=3$ statement,
as in \textsf{BS4-k}, and the suffix \textsf{-C} the $C$-fold counterpart,
in which $C$ disjoint $3$-collisions are asked for instead of one.  Two
further marks record a restriction of the \emph{algorithms} rather than of
the conclusion, and every statement carrying one is a theorem: index~$0$,
as in \textsf{BS0}, and a prime, as in \textsf{RS5'}, which is the
conclusion of \textsf{RS5} proved for a subclass.  Codes are also
\LaTeX{} labels, so \textsf{RS4} is \cref{RS4}.

The eight classes, with the state of play at $k=3$:

\begin{center}
\small
\begin{tabular}{@{}lllll@{}}
\toprule
class & memory constraint & mode & best algorithm & statements \\
\midrule
\cls{US} & none                  & seq. & $\Theta(N^{2/3})$ & settled,~\ref{F1} \\
\cls{UP} & none                  & par. & $T=\Theta(N^{2/3})$, $r=1$ & settled, \cref{p:onerou} \\
\cls{RS} & $S$ registers of points & seq. & $\tilde O(N/S)$ & \textsf{RS0}--\textsf{RS5}, \textsf{RS5'} \\
\cls{RP} & the same              & par. & $\tilde O(N/S)$ & none, \cref{p:costfree} \\
\cls{BS} & $S = \log_{2} W$ bits & seq. & $\tilde O(N/S)$ & \textsf{BS0}--\textsf{BS4}, \textsf{BS4-C} \\
\cls{BP} & the same              & par. & $\min\{N,\tilde\Theta(N^{2}/q^{2})\}$ & \textsf{BP0}--\textsf{BP3} \\
\cls{GS} & none, $\cmcs$ charged & seq. & $\tilde\Theta(N)$ & \textsf{GS1}--\textsf{GS4}, \textsf{GS3-C} \\
\cls{GP} & none, $\cmcs$ charged & par. & $\tilde\Theta(\max\{N^{2/3},N/q\})$ & \textsf{GP0}--\textsf{GP3} \\
\bottomrule
\end{tabular}
\end{center}

\medskip\noindent
In every constrained class, the best lower bound proved for \emph{all}
algorithms is the query bound of~\ref{F1}, which does not mention memory at
all.  Memory-sensitive bounds are proved only for restricted subclasses:
non-adaptive machines in \cref{RS0}, lifting ones in \cref{RS5'}, and
oblivious algorithms in \cref{BS0,BP0,GP0}.  Two classes are settled
because their algorithm meets the query bound, one is reducible, and five
carry ladders.

\begin{proposition}[one round is enough without a space bound]\label{p:onerou}
Querying $\Theta(N^{2/3})$ points at once and reading off a collision
succeeds with constant probability at $r = 1$.  So class \cls{UP} is
settled, at $T = \Theta(N^{2/3})$ and $r=1$, and nothing is open in it.
\end{proposition}

\noindent The one-round program of \cref{p:onerou} settles class \cls{UP}; the settled class \cls{US} is~\ref{F1}, and the reducible class \cls{RP} is \cref{sec:RP}.

\subsection{What is open, at a glance}

The rows below are ordered by how hard the lower bound is to prove, not by
how strong it is: the register-machine row is the easiest place to work and
the query-complexity row is already a theorem.  Every entry in the $k=2$
column is settled, and every entry to the right of it is open with the same
gap, namely the ratio between the best algorithm and the best proved lower
bound.

\begin{center}
\small
\begin{tabular}{@{}llll@{}}
\toprule
constraint & $k=2$ & $k=3$ & $k \geq 4$ \\
\midrule
none: query complexity
  & $\Theta(N^{1/2})$ \checkmark
  & $\Theta(N^{2/3})$ \checkmark
  & $\Theta(N^{(k-1)/k})$ \checkmark \\
register machine, $S = O(1)$
  & free \checkmark
  & gap $N^{1/3}$, \textsf{RS0}--\textsf{RS5}
  & gap $N^{1/k}$, \textsf{RS2-k}, \textsf{RS3-k} \\
branching program, $S = O(\log N)$
  & free \checkmark
  & gap $N^{1/3}$, \textsf{BS0}--\textsf{BS4}
  & gap $N^{1/k}$, \textsf{BS4-k} \\
cumulative memory, any $S$
  & $\tilde\Theta(N^{1/2})$ \checkmark
  & gap $N^{1/3}$, \textsf{GS1}--\textsf{GS4}
  & gap $N^{1/k}$, \textsf{GS3-k} \\
\addlinespace[2pt]
parallel, cumulative memory
  & free \checkmark
  & \textbf{open}, \textsf{GP0}--\textsf{GP3}
  & \textbf{open}, \textsf{GP3-k} \\
parallel, branching program
  & free \checkmark
  & \textbf{open}, \textsf{BP0}--\textsf{BP3}
  & \textbf{open}, \textsf{BP3-k} \\
\bottomrule
\end{tabular}
\end{center}

\medskip\noindent
Three readings.  \emph{The $k=2$ column is settled everywhere} because
rho attains the query bound with $O(\log N)$ bits, which is
Fact~\ref{F3}; that is why no conjecture here can be proved by an
argument indifferent to $k$.  \emph{The gap is $N^{1/k}$}, so it is widest
at $k=3$ and closes as $k$ grows, since the table-free recursion stays at
$\tilde O(N)$ while the query threshold rises.  And \emph{the three
constrained rows have identical known bounds}: they differ not in what is
proved but in how much room a proof has, which is the subject of
\cref{r:leverage}.

One memory-bounded cell is settled, and it is worth isolating because it is
the only one: for $C$ collision \emph{pairs} at $k=2$, Dinur~\cite{Din20}
proved $T^{2}S = \tilde\Theta(C^{2}N)$ in the branching-program model, in
the regime $S = \tilde O(C)$.  Every other proved lower bound in the table
is the query bound.

Filling in the $k=3$ column across the whole range of $S$, with the
Joux--Lucks algorithm as the upper bound:

\begin{center}
\small
\begin{tabular}{@{}lllll@{}}
\toprule
$S$ (cells) & best algorithm & proved lower bound & conjectured & status \\
\midrule
$O(1)$          & $\tilde O(N)$       & $\tilde\Omega(N^{2/3})$ & $\tilde\Theta(N)$       & open \\
$N^{1/6}$       & $\tilde O(N^{5/6})$ & $\tilde\Omega(N^{2/3})$ & $\tilde\Theta(N^{5/6})$ & open \\
$N^{1/4}$       & $\tilde O(N^{3/4})$ & $\tilde\Omega(N^{2/3})$ & $\tilde\Theta(N^{3/4})$ & open \\
$\geq N^{1/3}$  & $\tilde O(N^{2/3})$ & $\Omega(N^{2/3})$       & ---                     & settled \\
\bottomrule
\end{tabular}
\end{center}

\medskip\noindent
The proved column is constant because the only available bound is the
$S$-free one; the conjectured column is the algorithm, which is what
\cref{RS4} asserts to be optimal.  At $S = N^{1/3}$ the algorithm
meets the query bound and the row closes, so all of the content lies to the
left of that.

% ---------------------------------------------------------------------
% ---------------------------------------------------------------------
\section{Register machines}\label{sec:R}
% ---------------------------------------------------------------------

The register machine is the most tractable of the three models, for the
reason \cref{r:leverage} makes precise: its state holds points reached by
applying $f$, so one may reason in the functional graph of $f$ about what
the machine knows.  It is also, by \cref{app:classes}, an abstract model of
computation in Maurer's sense, the generic group model being the familiar
instance.  We treat the sequential machine first and in full; the parallel
machine, \cref{sec:RP}, reduces to it.

\subsection{Sequential: class \cls{RS}}\label{sec:reg}

\begin{cbox}
\textbf{Model.}  Fix $S$, $\beta$, and a sequence $w_1, w_2, \dots \in \NN$
chosen before $f$ is drawn.  A game holds a register file $Z[1..S]$ over
$\NN \cup \{\bot\}$, a control state $b \in \{0,1\}^{\beta}$, a head $hd$
into the sequence, and a cost counter; the algorithm holds nothing and acts
only through
\[
  \begin{array}{ll}
  \Advance(i,j)\colon & Z[j] \gets f(Z[i]) \quad\text{(charged)},\\
  \Next(j)\colon      & Z[j] \gets w_{hd},\quad hd \gets hd+1
                        \quad\text{(charged)},\\
  \Copy(i,j)\colon    & Z[j] \gets Z[i],\\
  \Update(b')\colon   & b \gets b' ,
  \end{array}
\]
halting by naming three register indices.  The algorithm is
\emph{stateless}: one function, fixed before $f$ is drawn, from the pair
$(b, v)$ to its next call, where $v \in \NN \cup \{\bot\}$ is the value
written by the operation it last issued, namely $f(Z[i])$ after
$\Advance(i,j)$, the cell $w_{hd}$ after $\Next(j)$, and $\bot$ otherwise.
Its \emph{cost} $T$ is the number of $\Advance$ and $\Next$ calls, and its
\emph{space} is $S$, counted in registers, each holding one point.  It
\emph{succeeds} if the three named registers hold a $3$-collision.
\Cref{app:models} gives the same model as code, and \cref{r:feedback} says
why the choice of $v$ as the report, rather than something weaker, is what
makes the algorithms of~\cite{JL09,OW99} legal in this class.
\end{cbox}

\medskip\noindent
Two clauses carry the model.  A register is written only by the three
charged or copying operations, so a point enters a register only as $f$ of
a point already present or as the next cell of the fixed sequence; and
$\Next$ is charged with a one-way head, so reaching the $m$-th cell costs
$m$.  Were $\Next$ free the machine could hold a desired index in its
control state, advance to it at no cost and query a point of its own
devising, which would collapse the model.  Everything else is free.
\Cref{app:models} gives the same model as code, and \cref{r:feedback} says
why the value written by the previous operation, rather than something
weaker, is the right report.

\begin{remark}[Where the leverage is]\label{r:leverage}
Class \cls{RS} is easier than class \cls{BS} for a structural reason.  In
a branching program the $S$ bits of state may encode anything, so one
cannot say what the program knows.  In a register machine the contents of
the state are points reached by applying $f$, so one may reason in the
functional graph of $f$: a $3$-collision is a vertex with three discovered
preimages, a single trajectory supplies two of them free at its own
cycle entry, and no trajectory supplies a third.  Quantitatively, for
trajectories of length $\ell$ a triple meets at a \emph{common} vertex with
probability $\approx \ell^3/N^2$ but merely fuses pairwise with probability
$\approx \ell^4/N^2$, a factor $\ell$ more often; the content of
\cref{RS3,RS4} is that this factor cannot be evaded.  Dinur's proof of
$T^2S = \tilde\Theta(C^2N)$ for finding $C$ collision
\emph{pairs}~\cite{Din20} is in the branching-program model and is the
technique to try for \cref{BS4}; he needed many pairs because for a single
pair, by~\ref{F3}, there is nothing to prove.
\end{remark}

\paragraph{The three-collision conjectures, $k=3$.}
The proved base is a counting bound and its lifting-machine strengthening;
the five conjectures then run from the weakest separation to the whole
advantage curve, which is also, by \cref{sec:route}, their order of
difficulty.

\begin{cthm}{RS0}[non-adaptive machines fail]\label{RS0}
Call a register machine \emph{non-adaptive} if its sequence of operations
and its halting time are fixed in advance, so that only the three indices
named on halting may depend on the register contents.  Then for every $S$,
every $\beta$ and every cost $T$, a non-adaptive machine succeeds with
probability $O(S^{3}/N^{2})$, which is $O(N^{-2})$ at $S = O(1)$.
\end{cthm}

\Cref{RS0} is why adaptivity is needed even to run the $\tilde O(N)$
algorithm of~\ref{F2}, which must stop when it finds its match.  It does not
separate the collision orders: the same count gives $O(S^{2}/N)$ at $k=2$,
where rho is likewise adaptive.

\begin{cthm}{RS5'}[lifting machines]\label{RS5'}
Call a successful run \emph{lifting} if, at the instant the last of the
three named points is first written into a register, the other two are
already present and their common image has been computed; that is, the run
finds a $2$-collision and then lifts it.  Call a machine lifting if all its
successful runs are.  A lifting register machine with $S$ registers and
cost $T$ succeeds with probability $O(TS/N)$, so \cref{RS5}, and hence
\cref{RS4}, hold for lifting machines.
\end{cthm}

\begin{proof}
Say a point is \emph{revealed} once the run has computed its image, and call
a value $y$ \emph{doubly covered} at a given time if the file then holds two
distinct revealed points with image $y$.  At most $\lfloor S/2 \rfloor$
values are doubly covered at any time, since each consumes two registers.
In a lifting run the last named point $x_3$ has $f(x_3) = y$ for some $y$
doubly covered just before $x_3$ is written.  Each charged operation reveals
the image of at most one point not revealed before, and conditioned on the
history that image is uniform on $\NN$, so it lands on a doubly covered
value with probability at most $S/2N$; if $x_3$ is never revealed its image
is uniform and independent as well, contributing $S/2N$ once.  A union bound
over the $T$ charged operations and the final naming gives $O(TS/N)$.
\end{proof}

\Cref{RS5'} covers every algorithm in the literature, including the
preimage-hunting machine that \cref{r:barrier} records as the obstruction in
class \cls{BS}: it is lifting with a single doubly covered value.  What
remains of the conjectures below is the runs that are \emph{not} lifting,
where one structure delivers all three preimages together, as a single
trajectory delivers two at its cycle entry by~\ref{F3}.

\medskip\noindent
The conjectures, weakest first.

\begin{conj}{RS1}[simplest unresolved]\label{RS1}
For every constant $S$ and every $\beta = O(\log N)$, a register machine
succeeding with probability $1-o(1)$ over the choice of $f$ has cost
$T = \omega\bigl(N^{2/3}\bigr)$.
\end{conj}
\begin{conj}{RS2}[qualitative]\label{RS2}
For every constant $S$ and every $\beta = O(\log N)$, a register machine
succeeding with probability at least $\tfrac12$ over the choice of $f$ has
cost $T = \omega\bigl(N^{2/3}\bigr)$.
\end{conj}
\begin{conj}{RS3}[memoryless]\label{RS3}
Under the hypotheses of \cref{RS2}, $T = N^{1-o(1)}$.
\end{conj}
\begin{conj}{RS4}[the whole curve]\label{RS4}
For every $S = S(N)$ and every $\beta = O(S\log N)$, a register machine
succeeding with probability at least $\tfrac12$ has cost
$T = \tilde\Omega\bigl(\max\{N^{2/3},\, N/S\}\bigr)$.
\end{conj}
\begin{conj}{RS5}[advantage at every cost]\label{RS5}
For every $S = S(N)$, every $\beta = O(S\log N)$ and every $T$, a register
machine of cost $T$ succeeds with probability $\tilde O(TS/N)$.
\end{conj}

\Cref{RS1} is the natural first target: a success probability of $1-o(1)$
lets one argue about typical executions rather than tail events, and the
qualitative form asks for no exponent.  \Cref{RS2} asks the same at constant
probability, \cref{RS3} pins the exponent and is tight by~\ref{F2}, and
\cref{RS4} replaces the point $S=O(1)$ by the whole curve and is tight by
Joux--Lucks; at $S \geq N^{1/3}$ it degenerates to~\ref{F1}, so all its
content lies below that.  \Cref{RS5} is the strongest: with the
unconditional $O\bigl((T^{3}+1)N^{-2}\bigr)$ of~\ref{F1} it bounds the
advantage by $\tilde O\bigl(\min\{TS/N,\ T^{3}/N^{2}\}\bigr)$, of which
\cref{RS4} is the constant-probability shadow.

\paragraph{General $k$.}
The counterparts replace $N^{2/3}$ by $N^{(k-1)/k}$ throughout; two are
worth stating, together with the whole-curve form, whose upper bound is now
attained by~\cite{GNY25}.

\begin{conj}{RS2-k}[qualitative, general $k$]\label{RS2-k}
For every constant $k \geq 3$, every constant $S$ and every
$\beta = O(\log N)$, a register machine succeeding with probability at
least $\tfrac12$ has cost $T = \omega\bigl(N^{(k-1)/k}\bigr)$.
\end{conj}
\begin{conj}{RS3-k}[memoryless, general $k$]\label{RS3-k}
Under the hypotheses of \cref{RS2-k}, $T = N^{1-o(1)}$.
\end{conj}
\begin{conj}{RS4-k}[the register curve, general $k$]\label{RS4-k}
For every constant $k \geq 3$, every $S = S(N)$ and every
$\beta = O(S\log N)$, a register machine succeeding with probability at
least $\tfrac12$ has cost
$T = \tilde\Omega\bigl(\max\{N^{(k-1)/k},\ N S^{-1/(k-2)}\}\bigr)$.
\end{conj}

Two cautions attach to the tightness of \cref{RS4-k}.  The upper bound is
specific to a domain and codomain of equal size: for $f\colon[N]\to[M]$ with
$M>N$ the curve breaks into two branches at $S=\tilde\Theta(N^{(k-2)/k})$,
and only the lower branch carries the exponent $1/(k-2)$~\cite{GNY25}.  And
the construction re-randomises its chains with auxiliary hash functions,
which class \cls{RS} does not supply; at $M=N$ these are replaceable by
fresh start points in the functional graph of $f$, as in parallel collision
search~\cite{OW99}, but the replacement has to be carried out before
\cref{RS4-k} may be called tight in this class.  The slack $N^{1/k}$ shrinks
as $k$ grows, so $k=3$ remains the easiest instance rather than merely the
first.

\subsection{Parallel: class \cls{RP}}\label{sec:RP}

For register machines parallelism is free, so class \cls{RP} carries no
conjecture of its own: it reduces to \cls{RS}.  A round issues $q$
operations on a file of $S$ cells.

\begin{proposition}[width is paid for in cells]\label{p:qcap}
The sources of the query operations in a round lie among the $S$ cells, so
a round issues at most $S$ distinct queries and one may assume
$q \leq S$.
\end{proposition}
\begin{proposition}[cost is unaffected by width]\label{p:costfree}
Running the $q$ operations of each round one at a time turns a $q$-parallel
register machine into a sequential one of the same cost and success
probability.  Hence \cref{RS4} covers every $q$, and class \cls{RP} carries
no statement of its own.
\end{proposition}
\begin{corollary}[depth]\label{cor:depth}
If \cref{RS4} holds then every $q$-parallel register machine succeeding
with probability at least $\tfrac12$ runs for
$r = \tilde\Omega\bigl(\max\{N^{2/3}/S,\; N/S^{2}\}\bigr)$ rounds.
\end{corollary}

\begin{proof}
Combine \cref{p:costfree} with $r=T/q$ and $q \leq S$ from \cref{p:qcap}.
\end{proof}

So for register machines parallelism buys latency and nothing else, and the
only quantity that genuinely depends on the width is cumulative memory,
class \cls{GP} of \cref{sec:GP}.

% ---------------------------------------------------------------------
\section{Branching programs}\label{sec:B}
% ---------------------------------------------------------------------

Dropping the restriction on \emph{which} point may be queried gives the
classical time--space model of Borodin and Cook~\cite{BC82}.  The state may
now encode anything, so one can no longer say what the program knows, which
is the source of the short-output barrier of \cref{r:barrier} and the reason
this model is harder to work in than the last.  Here the two modes genuinely
differ: \cref{r:asym} shows that a wide round changes the achievable bound,
so \cref{sec:BP} is not a reduction.

\subsection{Sequential: class \cls{BS}}\label{sec:bp}

\begin{cbox}
\textbf{Model.}  A \emph{branching program} of length $T$ and width $W$ is
a DAG on levels $0,\dots,T$ with at most $W$ nodes per level, in which
every node below level $T$ carries a label $x \in \NN$ and has one
outgoing edge for each $y \in \NN$, and every node at level $T$ carries a
label in $\NN^{3}$.  It is run by starting at the source, querying the
label of the current node, following the edge labelled by the answer, and
reading the label of the sink reached; it \emph{succeeds} if that label is
a $3$-collision.  Its \emph{space} is $S = \log_2 W$, counted in bits.  It
is \emph{oblivious} if all nodes at a level carry the same label, and
\emph{repetition-free} if no point labels two levels.  \Cref{app:models}
gives the same model as code, where these two restrictions become
statements about the type of one function.
\end{cbox}

\paragraph{The three-collision conjectures, $k=3$.}
The proved base is the oblivious curve, from the streaming switching lemma.
The genuinely open statements begin at \cref{BS1}; but the easiest thing to
attack in this class is not the weakest of them, it is the $C$-fold
statement below, the only one existing technique can address.

\begin{cthm}{BS0}[oblivious programs]\label{BS0}
Every repetition-free oblivious branching program succeeding with
probability at least $\tfrac12$ satisfies
$T \cdot \bigl(S + \lceil \log_2 N \rceil\bigr) = \tilde\Omega(N)$; in
particular $T = \tilde\Omega(N)$ at width $N^{O(1)}$.
\end{cthm}

\begin{proof}
The schedule $w_1,\dots,w_T$ is fixed in advance and its points are
distinct, so the answers arrive as a stream of $T$ elements drawn from $\NN$
uniformly with replacement.  Drawn instead without replacement, that is from
a random permutation, the stream would contain no repeated value and a
fortiori no $3$-collision.  Both distributions are invariant under permuting
positions, so the stream may be taken in schedule order.  Build a two-pass
distinguisher using $S+O(\log N)$ bits: the first pass runs the program and
retains the three points it names, the second reads their values and accepts
iff the three agree.  It never accepts a permutation stream, and it accepts
a function stream whenever the program succeeds, except for the $O(N^{-2})$
of~\ref{F1}.  The multi-pass streaming switching lemma
\cite[Thm.~5]{Din20b} bounds the advantage of a $P$-pass $S'$-bit
distinguisher on a stream of $T\leq N/3$ elements by $O\bigl(\log T\cdot
T(PS'+\log N)/N\bigr)$; with $P=2$ and $S'=S+O(\log N)$, an advantage of
$\tfrac12-O(N^{-2})$ forces $T(S+\lceil\log_2 N\rceil)=\tilde\Omega(N)$.
For $T>N/3$ the conclusion holds outright.
\end{proof}

At $k=3$, \cref{BS0} is tight only at $S=\mathrm{polylog}$; above that an
oblivious program can accumulate $P\leq S$ collision pairs at cost
$\tilde\Theta(PN/S)$ and hunt a third preimage in $\tilde\Theta(N/P)$, which
at $P=\sqrt S$ gives $T=\tilde\Theta(N/\sqrt S)$: a factor $\sqrt S$
separates \cref{BS0} from the truth in its own class, and closing it is the
smallest open quantity in this note.

\medskip\noindent
The genuinely open conjectures, in order of logical strength.

\begin{conj}{BS1}[weakest rung]\label{BS1}
Every branching program of width $N^{O(1)}$ succeeding with probability
$1-o(1)$ has length $T = \omega\bigl(N^{2/3}\bigr)$.
\end{conj}
\begin{conj}{BS2}\label{BS2}
The same at success probability at least $\tfrac12$.
\end{conj}
\begin{conj}{BS3}[polynomial width]\label{BS3}
Every branching program of width $N^{O(1)}$ succeeding with probability at
least $\tfrac12$ has length $T = N^{1-o(1)}$.
\end{conj}
\begin{conj}{BS4}[the whole curve]\label{BS4}
Every branching program succeeding with probability at least $\tfrac12$
over the choice of $f$ satisfies
$T \cdot S = \tilde\Omega(N)$ and $T = \tilde\Omega\bigl(N^{2/3}\bigr)$.
\end{conj}

The second clause of \cref{BS4} is~\ref{F1}; all the content is in
$T\cdot S$.  \Cref{BS3} is its polynomial-width case and reads simply: a
branching program of polynomial width needs length $N^{1-o(1)}$.
Randomisation need not be allowed for, since the input is already random.

\begin{remark}[Shapes]\label{r:shapes}
Comparing \cref{BS4} with what is proved is instructive.  Dinur's bound is
$T^2 S$; the conjecture here is $T S$.  The two agree at $k=2$ with one
pair, where $S = \tilde O(1)$ is inside his regime and $T^2S =
\tilde\Omega(N)$ gives $T = \tilde\Omega(\sqrt N)$, which is~\ref{F3}.
They differ in that the conjectured $k=3$ bound is linear in $S$ rather
than quadratic in $T$, which is exactly why $S = O(\log N)$ already forces
$T = \tilde\Omega(N)$ at $k=3$ and forces nothing at $k=2$.
\end{remark}
But none of \cref{BS1,BS2,BS3,BS4} is where to begin: by \cref{r:barrier}
each stands behind the same short-output wall, so the weakest is no easier
than the strongest.  The way in is to ask for many collisions.


\paragraph{The attackable case: many collisions.}\label{sec:cfold}
The Borodin--Cook method that \cref{BS4} would need cuts time into blocks and pays
a union bound over the $2^S$ nodes at a block boundary, affordable only if
producing a block's quota of outputs is improbable to a degree exponential
in the number of outputs.  A target of one triple offers no quota to be
improbable about; $C$ of them restore one.

\begin{conj}{BS4-C}[$C$ disjoint triples]\label{BS4-C}
For every $C \geq 1$, every branching program that outputs $C$ pairwise
disjoint $3$-collisions with probability at least $\tfrac12$ satisfies
\[
  T \;=\; \tilde\Omega\bigl(\max\{C^{1/3}N^{2/3},\ CN/S\}\bigr).
\]
\end{conj}

At $C=1$ this is \cref{BS4}.  It is tight: \cite{GNY25} find $C$ triples in
$\tilde O(CN/S)$ for every $S\leq\tilde O(C^{2/3}N^{1/3})$, and there the
query threshold $\Theta((CN^2)^{1/3})$ is reached, so the two clauses cross
where the two regimes meet.  The reason to state it is that it is the first
statement of class \cls{BS} existing technique can address at all: the block
union bound needs $\tilde\Omega(S)$ outputs per block, so the method is
available once $C=\tilde\Omega(S)$, and $C=S=\mathrm{polylog}$ lies inside
that while $C=1$ does not.  The companion note \cite{Note2} works this case
and shows that segmentation delivers the weaker
$T\cdot S^{2/3}=\tilde\Omega(CN^{2/3})$, short of \cref{BS4-C} by the same
$N^{1/3}$ that separates query bound from algorithm at $C=1$; what is
missing is an argument that a short program cannot \emph{hold} many pairs,
as opposed to \emph{find} them.

\paragraph{General $k$.}

\begin{conj}{BS4-k}[branching programs, general $k$]\label{BS4-k}
For every constant $k \geq 3$, every branching program succeeding with
probability at least $\tfrac12$ satisfies
\[
  T \cdot S^{1/(k-2)} \;=\; \tilde\Omega(N)
  \qquad\text{and}\qquad
  T \;=\; \tilde\Omega\bigl(N^{(k-1)/k}\bigr).
\]
\end{conj}

\Cref{BS4-k} is tight at both attested endpoints for every $k$, and
throughout for $k=3$; for $k\geq4$ its interior rests on the power law
$T=\tilde\Theta(NS^{-1/(k-2)})$ and should be read as the natural guess
rather than as calibrated against an algorithm.  The remaining rungs have
the counterparts got by replacing $N^{2/3}$ with $N^{(k-1)/k}$, and are not
restated.

\subsection{Parallel: class \cls{BP}}\label{sec:BP}

Here a node carries a $q$-tuple and branches on the whole answer tuple, so
it may name any number of points regardless of the width; width is free, and
that changes the answer.

\begin{proposition}[width is free, and lowers the cost]\label{p:parbp}
There is a one-round parallel branching program of width $O(N^{2})$ that
queries $T = \Theta(N^{2/3})$ points and succeeds with constant
probability: query $T$ fixed points, keep one sink per candidate triple, and
route each answer tuple to a sink that collides under it.
\end{proposition}
\begin{remark}[Where the asymmetry comes from]\label{r:asym}
A wide round lets a program \emph{select} its output from the whole answer
batch, and selecting is cheap.  A register machine must instead \emph{hold}
its output in cells, and by \cref{p:qcap} the cell count caps how much of a
batch it can hold, so buying width costs memory.  This is why
\cref{p:costfree} has no analogue in class \cls{BP}, and it is the third
reason, after tractability and the leverage of \cref{r:leverage}, to prefer
the register model.
\end{remark}

\paragraph{What the parallel programs achieve.}
For a program of polynomial width, two strategies are available, each with
$O(\log N)$ bits between rounds.  One holds a $2$-collision and hunts $q$
fresh points per round, costing $T=\Theta(N)$ at any $q$.  The other detects
a $3$-collision inside a single batch, which routing permits once the width
exceeds $\binom{q}{3}$; the probability is $\approx q^3/N^2$ per round, so
$r\approx N^2/q^3$ and $T=\tilde\Theta(N^2/q^2)$.  Taking the better of the
two, the crossover is at $q=\sqrt N$.

\paragraph{The three-collision conjectures, $k=3$.}
The proved base is again the oblivious curve; the conjectures then run from
one width to the whole range.

\begin{cthm}{BP0}[oblivious parallel programs]\label{BP0}
Every repetition-free oblivious parallel branching program of space $S$ and
query width $q$ succeeding with probability at least $\tfrac12$ satisfies
$T \cdot \bigl(S + q\lceil\log_2 N\rceil\bigr) = \tilde\Omega(N)$; in
particular $T = \tilde\Omega(N/q)$ at width $N^{O(1)}$.
\end{cthm}

\begin{proof}
As in \cref{BS0}.  A round branches on its whole answer tuple, so read one
element at a time the program is a streaming algorithm carrying
$S+q\lceil\log_2N\rceil$ bits across the interior of a round; the verifying
second pass needs $O(\log N)$ more.
\end{proof}

\Cref{BP0} settles the oblivious case and falls short of \cref{BP3} below by
a factor $q$ even there.  The conjectures, weakest first.

\begin{conj}{BP1}[the crossing width]\label{BP1}
Every parallel branching program of width $N^{O(1)}$ and query width
$q = \sqrt N$ succeeding with probability at least $\tfrac12$ satisfies
$T = \tilde\Omega(N)$, that is $r = \tilde\Omega(\sqrt N)$ rounds.
\end{conj}
\begin{conj}{BP2}[the hold-and-hunt range]\label{BP2}
The same for every $q \leq \sqrt N$, where hold-and-hunt is the better
strategy and the target is $\tilde\Omega(N)$ throughout.
\end{conj}
\begin{conj}{BP3}[every width]\label{BP3}
For every $q$, every parallel branching program of width $N^{O(1)}$
succeeding with probability at least $\tfrac12$ satisfies
$T = \tilde\Omega\bigl(\min\{\, N,\; N^{2}/q^{2} \,\}\bigr)$.
\end{conj}

\Cref{BP1} is the logically weakest but not the easiest instance: $q=\sqrt
N$ is where the two strategies cross and the curve is hardest, whereas on
either side one strategy is comfortably better and an argument has more
room, which is \cref{r:parbase}.  \Cref{BP3} at $q=1$ is the
polynomial-width case of \cref{BS4}, so it is stronger.  Beyond polynomial
width the shape is unclear, since a larger width both permits a table and
permits naming more triples per batch; we state no conjecture there, and
note that the naive $T\cdot(S+q\lceil\log_2N\rceil)=\tilde\Omega(N)$ is
proved only in the oblivious case (\cref{BP0}) and is loose regardless, by
\cref{r:notconj}.

\begin{remark}[What is not conjectured]\label{r:notconj}
For branching programs of width beyond polynomial the shape is unclear,
because the two strategies interact: a larger width both permits a table,
as in the sequential case, and permits naming more triples per batch.  We
state no conjecture there.  Note also that the bound
$T \cdot (S + q\lceil\log_2 N\rceil) = \tilde\Omega(N)$, which is the first
thing one writes down, is established only in the oblivious case, where it
is \cref{BP0}; for adaptive programs it follows from \cref{BS4} and from
nothing else proved here.  Either way it is loose: at $q = N^{1/2}$ it
gives $T = \tilde\Omega(N^{1/2})$ against a truth of $\tilde\Theta(N)$, so
even \cref{BP0} leaves a factor $\sqrt N$ inside its own class at the
crossing width.
\end{remark}

\begin{remark}[Which base rung to attack]\label{r:parbase}
\Cref{GP1} is stated at $q = N^{1/6}$ merely as a clean interior point; any
$1 \ll q \ll N^{1/3}$ would do, the endpoints being respectively
\cref{GS3} and the query charge.  \Cref{BP1} is different: it is the
logically weakest statement of its class, being a single width, but it is
not the easiest instance, since $q = \sqrt N$ is where the two strategies
of \cref{sec:BP} cross and the curve is hardest.  On either side of that
width one strategy is comfortably better than the other, so an argument
that need only rule out near-optimal programs has more room there.
\end{remark}

\paragraph{General $k$.}

\begin{conj}{BP3-k}[parallel branching programs, general $k$]\label{BP3-k}
For every constant $k \geq 3$ and every $q$, every parallel branching
program of width $N^{O(1)}$ succeeding with probability at least $\tfrac12$
satisfies
$T = \tilde\Omega\bigl(\min\{\, N,\; N^{k-1}/q^{k-1} \,\}\bigr)$.
\end{conj}

The batch-detection strategy becomes $r\approx N^{k-1}/q^k$ and hence
$T=\tilde\Theta(N^{k-1}/q^{k-1})$, so the two strategies cross at
$q=N^{(k-2)/(k-1)}$, which is $\sqrt N$ at $k=3$.

% ---------------------------------------------------------------------
\section{Cumulative memory}\label{sec:G}
% ---------------------------------------------------------------------

The models above bound the memory available at each step; cumulative memory
integrates memory over time, which is the measure under which
memory-hardness is normally stated~\cite{AS15}.  By \cref{r:transfer} a
Borodin--Cook proof in class \cls{BS} transfers here with only a constant or
logarithmic loss, so class \cls{GS} is strictly stronger than \cls{BS} and
no harder to work in; the two are one stage of \cref{sec:route}, not two.

\subsection{Sequential: class \cls{GS}}\label{sec:cmc}

\begin{cbox}
\textbf{Model.}  A sequential adversary is a sequence of functions
$\advers_i$ carrying a state $\sigma_i \in \{0,1\}^{*}$ from one query to
the next; after $T$ queries it outputs three points.  Its \emph{charged
cumulative memory} on a given run is
\[
  \cmcs \;=\; \sum_{i=1}^{T} |\sigma_i| \;+\; T\lceil \log_2 N \rceil ,
\]
its \emph{expected} charged cumulative memory $\eCMCs$ is the expectation
of $\cmcs$ over $f$ and the adversary's coins, and its \emph{worst-case}
charged cumulative memory $\CMCs$ is the supremum of $\cmcs$ over $f$ and
those coins; so $\eCMCs \leq \CMCs$.  Its \emph{peak space} is
$\max_i |\sigma_i|$, and a peak space of $S$ \emph{cells} means
$S\lceil\log_2 N\rceil$ bits, the distinction being absorbed by
$\tilde\Omega$.  The second term of $\cmcs$ charges the queries
themselves; without it the measure admits no lower bound, since an
adversary may issue all its queries in one batch while carrying no state.
\Cref{app:models} gives the measure as an accumulator inside the execution
game.
\end{cbox}

\paragraph{The three-collision conjectures, $k=3$.}
There is no proved base in this class: the conversion of \cref{r:transfer}
delivers one as soon as \cref{BS4-C} is proved, but nothing is proved
unconditionally beyond the query bound.  The conjectures run from one
interior peak space in the worst-case metric to the advantage form; and,
as in class \cls{BS}, the attackable statement is the $C$-fold one.

\begin{conj}{GS1}[one interior peak space, worst case]\label{GS1}
Every sequential adversary of peak space $S = N^{1/6}$ succeeding with
probability at least $\tfrac12$ satisfies $\CMCs = \tilde\Omega(N)$.
\end{conj}
\begin{conj}{GS2}[the same in the expected metric]\label{GS2}
Every sequential adversary of peak space $S = N^{1/6}$ succeeding with
probability at least $\tfrac12$ satisfies $\eCMCs = \tilde\Omega(N)$.
\end{conj}
\begin{conj}{GS3}[every peak space]\label{GS3}
Every sequential adversary succeeding with probability at least $\tfrac12$
satisfies $\eCMCs = \tilde\Omega(N)$.
\end{conj}
\begin{conj}{GS4}[advantage]\label{GS4}
For every $\epsilon \in (0,1]$, every sequential adversary succeeding with
probability at least $\epsilon$ satisfies $\eCMCs=\tilde\Omega(\epsilon N)$.
\end{conj}

\Cref{GS3} is tight by~\ref{F2} and by Joux--Lucks at every point of its
curve, since $T\cdot S=\tilde\Theta(N)$ there and the state is held for
essentially the whole run; that the whole family sits at one value of
$\cmcs$ is the reason to prefer this measure.  \Cref{GS1,GS2} must be stated
at an interior $S$: at $S=O(\log N)$ bits cumulative memory degenerates into
the query charge and the statement there is \cref{BS3}.  In the expected
metric the linear factor of \cref{GS4} is necessary, by \cref{r:eps}.

\begin{remark}[The factor $\epsilon$]\label{r:eps}
In the expected metric the linear factor of \cref{GS4} is necessary: an
adversary that runs the algorithm of~\ref{F2} with probability $\epsilon$
and halts at once otherwise succeeds with probability
$\Omega(\epsilon)$ and has $\eCMCs = \tilde\Theta(\epsilon N)$.  In the
worst-case metric that construction pays in full on the branch where it
runs, so no such factor appears and the analogue of \cref{GS4} is
$\CMCs = \tilde\Omega(N)$; the limit of any such statement is the guessing
floor at $\epsilon = \Theta(N^{-2})$, where naming three unqueried points
already succeeds.
\end{remark}

\begin{remark}[Why the base rungs sit at an interior peak space]\label{r:interior}
\Cref{GS1,GS2} must be stated at an interior $S$.  At $S = O(\log N)$ bits
cumulative memory degenerates into the query charge and the statement there
is \cref{BS3} rather than anything new.  At $S = N^{1/6}$ the target is
tight, since Joux--Lucks there has $T = N^{5/6}$ and
$\cmcs \approx TS = \tilde\Theta(N)$; any $1 \ll S \ll N^{1/3}$ would serve
equally.
\end{remark}

\begin{remark}[Between peak and cumulative]\label{r:sustained}
Peak space and cumulative memory are not the only two measures.  Sustained
space~\cite{ABP18}
counts the steps at which the space in use is at least $s$, and a statement
of the form \emph{every successful adversary spends $\tilde\Omega(N/s)$
steps with space at least $s$} implies \cref{GS3} already for one interior
$s$, since those steps alone contribute $\cmcs = \tilde\Omega(N)$.  It is
therefore strictly stronger than \cref{GS3}, and it is the natural statement
to attempt once \cref{GS3} is in hand: the transfer of \cref{r:transfer} is
stated for cumulative memory, but the block decomposition it rests on is
what a sustained bound needs too.
\end{remark}


\paragraph{The attackable case: many collisions.}

\begin{conj}{GS3-C}[the same in cumulative memory]\label{GS3-C}
For every $C \geq 1$, every sequential adversary that outputs $C$ pairwise
disjoint $3$-collisions with probability at least $\tfrac12$ satisfies
$\eCMCs = \tilde\Omega(CN)$.
\end{conj}

At $C=1$ this is \cref{GS3}, and by \cref{r:transfer} a proof of
\cref{BS4-C} along the expected lines carries it at no extra cost.  It is
tight by the same algorithm as \cref{BS4-C}.

\paragraph{General $k$.}

\begin{conj}{GS3-k}[cumulative memory, general $k$]\label{GS3-k}
For every constant $k \geq 3$, every sequential adversary succeeding with
probability at least $\tfrac12$ satisfies $\eCMCs = \tilde\Omega(N)$.
\end{conj}

\Cref{GS3-k} deserves a word: its target does \emph{not} depend on $k$.  It
is tight for every $k\geq3$, since the table-free recursion has
$\cmcs=\tilde\Theta(N)$ while the balanced Joux--Lucks point has
$\cmcs\approx TS=N^{(2k-3)/k}\geq N$, with equality exactly at $k=3$.  So
cumulative memory sees no difference between the collision orders above two,
and at $k=3$ alone the whole tradeoff curve sits at one value of it.

\subsection{Parallel: class \cls{GP}}\label{sec:GP}

A round issues $q$ queries, the length becomes the number of rounds $r$, and
$T=rq$.  This is the one class in which the achievable cumulative memory
genuinely depends on the width.

\paragraph{The three-collision conjectures, $k=3$.}
The proved base is the oblivious bound; the conjectures then run from one
width to the whole range.

\begin{cthm}{GP0}[oblivious parallel adversaries]\label{GP0}
Every repetition-free oblivious $q$-parallel adversary of peak space $S$
cells succeeding with probability at least $\tfrac12$ satisfies
$T \cdot (S + q) = \tilde\Omega(N)$, and hence
$\eCMCs = \tilde\Omega\bigl(N/(S+q)\bigr)$.  At $S = \tilde O(q)$ and
$q \leq N^{1/3}$ this is \cref{GP3}.
\end{cthm}

\begin{proof}
The first clause is \cref{BP0}, whose proof uses only the state carried
between rounds and the buffer held within one, and which applies at any
constant success probability.  For the second, $\cmcs\geq
T\lceil\log_2N\rceil$ pointwise, so an adversary with
$\eCMCs=o\bigl(N/(S+q)\bigr)$ has $\mathbb{E}[T]=o\bigl(N/(S+q)\bigr)$;
truncating it at $4\,\mathbb{E}[T]$ queries leaves success probability at
least $\tfrac14$ by Markov's inequality, contradicting the first clause.
\end{proof}

The conjectures, weakest first.

\begin{conj}{GP1}[one interior width, worst case]\label{GP1}
Every $q$-parallel adversary of width $q = N^{1/6}$ succeeding with
probability at least $\tfrac12$ satisfies $\CMCs = \tilde\Omega(N^{5/6})$.
\end{conj}
\begin{conj}{GP2}[the same in the expected metric]\label{GP2}
Every $q$-parallel adversary of width $q = N^{1/6}$ succeeding with
probability at least $\tfrac12$ satisfies $\eCMCs = \tilde\Omega(N^{5/6})$.
\end{conj}
\begin{conj}{GP3}[every width]\label{GP3}
For every $q \geq 1$, every $q$-parallel adversary succeeding with
probability at least $\tfrac12$ satisfies
$\eCMCs = \tilde\Omega\bigl(\max\{N^{2/3},\, N/q\}\bigr)$; equivalently
$q \cdot \eCMCs = \tilde\Omega(N)$ for $q \leq N^{1/3}$.
\end{conj}

\Cref{GP3} at $q=1$ is \cref{GS3}, so it is stronger; it is met by
Joux--Lucks at width $q$, where $\cmcs\approx rS=N/q$, and at $q=N^{1/3}$ it
flattens onto the query charge.  \Cref{GP1} is stated at $q=N^{1/6}$ as a
clean interior point; any $1\ll q\ll N^{1/3}$ would do.

\paragraph{General $k$.}

\begin{conj}{GP3-k}[cumulative memory, parallel, general $k$]\label{GP3-k}
For every constant $k \geq 3$ and every $q \geq 1$, every $q$-parallel
adversary succeeding with probability at least $\tfrac12$ satisfies
$\eCMCs = \tilde\Omega\bigl(\max\{N^{(k-1)/k},\, N/q\}\bigr)$; equivalently
$q \cdot \eCMCs = \tilde\Omega(N)$ for $q \leq N^{1/k}$.
\end{conj}

The target of \cref{GP3-k} is again $k$-free away from the floor.
% ---------------------------------------------------------------------
\section{Relations across the classes}\label{sec:relations}
% ---------------------------------------------------------------------

The model sections presented each class in its own order of difficulty.
This section records the logical relations that cut across them: the ladders
within each class, the chain between classes, how both change with $k$, the
structural reason register machines are the easiest place to work and
branching programs the hardest, and what a single algorithm would have to do
to refute each statement.

\begin{proposition}[the ladders]\label{p:ladders}
Within each class the codes are ordered by implication:
\begin{center}
\begin{tabular}{@{}l@{}}
\ref{RS5} $\Rightarrow$ \ref{RS4} $\Rightarrow$ \ref{RS3} $\Rightarrow$
  \ref{RS2} $\Rightarrow$ \ref{RS1},\\
\ref{BS4-C} $\Rightarrow$ \ref{BS4} $\Rightarrow$ \ref{BS3} $\Rightarrow$
  \ref{BS2} $\Rightarrow$ \ref{BS1},\\
\ref{GS4} $\Rightarrow$ \ref{GS3} $\Rightarrow$ \ref{GS2} $\Rightarrow$
  \ref{GS1}, \quad and \quad \ref{GS3-C} $\Rightarrow$ \ref{GS3},\\
\ref{GP3} $\Rightarrow$ \ref{GP2} $\Rightarrow$ \ref{GP1},\\
\ref{BP3} $\Rightarrow$ \ref{BP2} $\Rightarrow$ \ref{BP1}.
\end{tabular}
\end{center}
The proved statements \cref{RS0,RS5',BS0,BP0,GP0} are not rungs of these
ladders.
\end{proposition}

\begin{proof}
Four moves generate every implication.  \emph{Weakening the hypothesis on
the algorithm}: demanding success probability $1-o(1)$ rather than $\tfrac12$
restricts the algorithms quantified over, giving \ref{RS2} $\Rightarrow$
\ref{RS1} and \ref{BS2} $\Rightarrow$ \ref{BS1}.  \emph{Weakening the
conclusion}: $N^{1-o(1)}=\omega(N^{2/3})$, giving \ref{RS3} $\Rightarrow$
\ref{RS2} and \ref{BS3} $\Rightarrow$ \ref{BS2}.  \emph{Specialising a curve
to one value of its parameter}: $S=O(1)$ in \cref{RS4} gives \cref{RS3};
$S=O(\log N)$ in \cref{BS4} gives \cref{BS3}; $S=N^{1/6}$ in \cref{GS3} gives
\cref{GS2}; $q=N^{1/6}$ in \cref{GP3} gives \cref{GP2}; $q\leq\sqrt N$ in
\cref{BP3} gives \cref{BP2} and $q=\sqrt N$ then \cref{BP1};
$\epsilon=\tfrac12$ in \cref{GS4} gives \cref{GS3}; $C=1$ in \cref{BS4-C}
gives \cref{BS4} and in \cref{GS3-C} gives \cref{GS3}; and success
probability $\tfrac12$ in \cref{RS5} gives the clause $T=\tilde\Omega(N/S)$
of \cref{RS4}, whose other clause is~\ref{F1}.  \emph{Passing from the
expected to the worst-case metric}: $\eCMCs\leq\CMCs$, giving \ref{GS2}
$\Rightarrow$ \ref{GS1} and \ref{GP2} $\Rightarrow$ \ref{GP1}.  The five
proved statements restrict the algorithms rather than the conclusion and so
sit beside the ladders rather than on them; \cref{sec:route} places them.
\end{proof}

\begin{proposition}[the chain across classes]\label{p:chain}
\ref{GS3} $\Rightarrow$ \ref{BS4} $\Rightarrow$ \ref{RS4}, and hence with
\cref{p:ladders} the sequential statements have \ref{GS4} at the top and
\ref{RS1} at the bottom.
\end{proposition}

\begin{proof}
For the first, $\cmcs\leq T(S+\lceil\log_2 N\rceil)$ pointwise for an
adversary of peak space $S$, so $\eCMCs=\tilde\Omega(N)$ forces
$T\cdot\max\{S,\log_2 N\}=\tilde\Omega(N)$.  For the second, a register
machine with $S$ cells and $\beta=O(S\log N)$ control bits is a branching
program of the same length and space $O(S\log N)$, its node at level $i$
being the pair $(Z,b)$; conversely a branching program need not be a
register machine, since its node may be labelled with any point.
\end{proof}

\begin{remark}[Peak does not bound cumulative]\label{r:peak}
\Cref{BS4} does not imply \cref{GS3}.  A bound on $T \cdot S$ constrains
the \emph{peak} state, and $\sum_i |\sigma_i|$ can be far below $TS$ when
the state is small in most rounds.  Any route from a branching-program
bound to a cumulative one therefore needs a further argument about the
profile $i \mapsto |\sigma_i|$, not just its maximum.  This is why
\cref{GS3} is listed as strictly stronger and not as a corollary.
\end{remark}
\begin{remark}[The technique transfers even though the statement does not]\label{r:transfer}
\Cref{r:peak} is about implication, not about difficulty.  Beame and
Kornerup~\cite{BK23} convert Borodin--Cook-style time--space proofs into
cumulative-memory bounds of the same strength, by choosing the block
boundaries at time steps of locally small space instead of at fixed ones,
and losing only a constant or logarithmic factor.  So a proof of
\cref{BS4-C} by that route delivers \cref{GS3-C}, and a proof of
\cref{BS4} by that route would deliver \cref{GS3}.  Class \cls{GS} is
strictly stronger than class \cls{BS} and no harder to work in, which is
why \cref{sec:route} treats their rungs as one stage and not two.  The step
is not speculative: applied to the estimate of Hamoudi and Magniez for $k$
disjoint collision pairs it yields the cumulative-memory bound
$\Omega(k^{3}N/T^{2})$ for quantum circuits~\cite{HM21,BK23}, which is the
collision-pair instance of exactly the move \cref{BS4-C} would need.
\end{remark}

\begin{remark}[The short-output barrier]\label{r:barrier}
Every statement of classes \cls{BS}, \cls{GS}, \cls{BP} and \cls{GP} asks
for one $3$-collision, an output of $O(\log N)$ bits, and that is the
regime in which no technique for general models is known to give anything:
the available tradeoffs need an output at least as large as the space, and
for short outputs the best bounds are logarithmic in the input
size~\cite{HHL26}.  Hao, Huang and Liu take the $3$-collision problem as
their starting point and cannot treat it.  The reduction one wants, from
finding one triple to finding many pairs, fails on machines that never hold
many pairs, their example being one that hunts preimages of a single value
in polylogarithmic space; they prove a tradeoff instead for a modified
problem in which a second oracle records the tuples an algorithm has
assembled.  Two consequences here.  That example is a constraint alongside
\ref{F3} which any proof must accommodate, and \cref{RS5'} accommodates it
by covering exactly the machines it describes.  And the statements of
\cref{sec:cfold} are the way past the barrier that leaves the problem
alone.  By contrast streaming models carry no such barrier, which is why
\ref{BS0}, \ref{BP0} and \ref{GP0} are proved; the preprocessing model
carries none either, but it answers a different question, as
\cref{r:proxies} records.  The quantum counterpart of the same wall is
where~\cite{HHL26} began, the time--space complexity of quantum collision
finding being on Aaronson's list of ten open challenges.
\end{remark}
\begin{remark}[Three neighbouring models that do not see the phenomenon]\label{r:proxies}
Each of the following questions has been settled, and none of them settles
anything here; it is worth being explicit about why, since all three are
sometimes offered as proxies.

\emph{Preprocessed advice.}  In the auxiliary-input random oracle model the
adversary receives $S$ bits about $f$ and then queries with unbounded
working memory.  Memory-sensitive bounds are proved there, and for
multi-collisions in Merkle--Damg\aa{}rd hashing the $m$-way advantage
differs from the $2$-way one by a factor $m$~\cite{ACDW20,Aks24}.  The
phenomenon of this note is invisible in that model, because advice is not
working memory: an adversary with polylogarithmic advice still has the
$\Theta(N^{2/3})$-query algorithm of~\ref{F1} available.

\emph{Deciding instead of searching.}  For a random $f \colon \NN \to \NN$ a
$3$-collision exists with overwhelming probability, so the decision problem
is vacuous.  In the non-trivial regime $M = \tilde\Theta(N^{3/2})$ for
$f \colon [N] \to [M]$ it is solvable in $\tilde O(N^{3/2}/\sqrt S)$, which
is the curve for $2$-collisions as well~\cite{GNY25}, so deciding does not
see the collision order either.

\emph{Comparison-based programs.}  For element distinctness, a single-output
problem, memory-sensitive bounds are proved: $TS = \Omega(n^{3/2}\sqrt{\log
n})$~\cite{BFMUW87}, improved to $TS = \Omega(n^{2-o(1)})$~\cite{Yao94}.
Both are for comparison-based programs on worst-case inputs.  The problem
here differs on three axes at once, being a search problem, on a random
input, against programs that may compute anything of what they have seen,
and each difference is one of the reasons \cref{r:barrier} bites.
\end{remark}
\begin{proposition}\label{p:fork}
For $k = 3$ the general-$k$ statements form the single chain of
\cref{p:chain}.  For $k \geq 4$ they do not:
\ref{BS4-k} $\Rightarrow$ \ref{RS3-k} $\Rightarrow$ \ref{RS2-k}, and
\ref{GS3-k} $\Rightarrow$ \ref{BS4-k} \emph{fails}.
\end{proposition}

\begin{proof}
The surviving implications are as in \cref{p:chain}, a register machine with
$S$ cells being a branching program of space $\tilde O(S)$.  For the failure,
$\cmcs\leq T(S+\lceil\log_2 N\rceil)$ yields only $T\cdot S=\tilde\Omega(N)$,
and since $S^{1/(k-2)}\leq S$ that is strictly weaker than the conclusion of
\cref{BS4-k} whenever $k\geq4$.
\end{proof}

So at $k=3$, where $1/(k-2)=1$, the cumulative and branching-program
statements coincide in strength and the sequential statements are linearly
ordered; from $k=4$ on they are incomparable.


\subsection{What would refute each}


Each statement is refuted by a single algorithm, and it is worth being
explicit about what would count.  Throughout $k=3$.

\begin{center}
\small
\begin{tabular}{@{}ll@{}}
\toprule
statement & refuted by \\
\midrule
\ref{RS2} & a register machine with $S = O(1)$ and cost $O(N^{2/3})$ \\
\ref{RS3} & the same with cost $N^{1-\delta}$ for a fixed $\delta > 0$ \\
\ref{RS4} & a register machine beating $T = \tilde O(N/S)$ for some $S$ \\
\ref{RS5} & any of the above, or an advantage $\omega(TS/N)$ at some $T$ \\
\ref{BS4} & any algorithm with $T \cdot S = o(N/\mathrm{polylog})$ \\
\ref{GS3} & any algorithm with $\eCMCs = o(N/\mathrm{polylog})$ \\
\ref{GP3} & a $q$-parallel adversary with $\eCMCs = o(\max\{N^{2/3},N/q\})$ \\
\ref{BP3} & a $q$-parallel program of width $N^{O(1)}$ with
            $T = o(\min\{N,N^{2}/q^{2}\})$ \\
\ref{BS4-C} & any algorithm with $T \cdot S = o(CN/\mathrm{polylog})$ for some $C$ \\
\ref{GS3-C} & any algorithm with $\eCMCs = o(CN/\mathrm{polylog})$ for some $C$ \\
\bottomrule
\end{tabular}
\end{center}
\section{Index of labelled statements}\label{sec:index}
% ---------------------------------------------------------------------

\begin{center}
\footnotesize
\begin{tabular}{@{}llp{0.44\textwidth}l@{}}
\toprule
code & kind & statement & status \\
\midrule
\multicolumn{4}{@{}l}{\emph{class \cls{RS}, register machine, sequential
  (\cref{sec:reg})}}\\
\ref{RS0} & Thm  & non-adaptive: success $O(S^{3}/N^{2})$ at every $T$ & proved \\
\ref{RS1} & Conj & $S=O(1)$, success $1-o(1)$: $T=\omega(N^{2/3})$ & open, base rung \\
\ref{RS2} & Conj & $S=O(1)$, success $\tfrac12$: $T=\omega(N^{2/3})$ & open \\
\ref{RS3} & Conj & $S=O(1)$, success $\tfrac12$: $T=N^{1-o(1)}$ & open \\
\ref{RS4} & Conj & every $S$: $T=\tilde\Omega(\max\{N^{2/3},N/S\})$ & open \\
\ref{RS5} & Conj & every $T$: advantage $\tilde O(TS/N)$ & open, top rung \\
\ref{RS5'} & Thm  & lifting machines: advantage $O(TS/N)$ & proved \\
\addlinespace[3pt]
\multicolumn{4}{@{}l}{\emph{class \cls{BS}, branching program, sequential
  (\cref{sec:bp})}}\\
\ref{BS0} & Thm  & repetition-free oblivious: $T\cdot(S+\log_2 N)=\tilde\Omega(N)$ & proved \\
\ref{BS1} & Conj & width $N^{O(1)}$, success $1-o(1)$: $T=\omega(N^{2/3})$ & open, weakest rung \\
\ref{BS2} & Conj & width $N^{O(1)}$, success $\tfrac12$: $T=\omega(N^{2/3})$ & open \\
\ref{BS3} & Conj & width $N^{O(1)}$: $T=N^{1-o(1)}$ & open \\
\ref{BS4} & Conj & every width: $T\cdot S=\tilde\Omega(N)$ & open, top rung \\
\ref{BS4-C} & Conj & $C$ disjoint triples: $T=\tilde\Omega(\max\{C^{1/3}N^{2/3},CN/S\})$ & open, reachable at $C=\tilde\Omega(S)$ \\
\addlinespace[3pt]
\multicolumn{4}{@{}l}{\emph{class \cls{GS}, general adversary, sequential
  (\cref{sec:cmc})}}\\
\ref{GS1} & Conj & peak $S=N^{1/6}$: $\CMCs=\tilde\Omega(N)$ & open, base rung \\
\ref{GS2} & Conj & peak $S=N^{1/6}$: $\eCMCs=\tilde\Omega(N)$ & open \\
\ref{GS3} & Conj & every peak space: $\eCMCs=\tilde\Omega(N)$ & open \\
\ref{GS4} & Conj & success $\epsilon$: $\eCMCs=\tilde\Omega(\epsilon N)$ & open, top rung \\
\ref{GS3-C} & Conj & $C$ disjoint triples: $\eCMCs=\tilde\Omega(CN)$ & open, follows \ref{BS4-C} \\
\addlinespace[3pt]
\multicolumn{4}{@{}l}{\emph{class \cls{GP}, general adversary, parallel
  (\cref{sec:GP})}}\\
\ref{GP0} & Thm  & repetition-free oblivious: $T\cdot(S+q)=\tilde\Omega(N)$ & proved \\
\ref{GP1} & Conj & $q=N^{1/6}$: $\CMCs=\tilde\Omega(N^{5/6})$ & open, weakest rung \\
\ref{GP2} & Conj & $q=N^{1/6}$: $\eCMCs=\tilde\Omega(N^{5/6})$ & open \\
\ref{GP3} & Conj & every $q$: $\eCMCs=\tilde\Omega(\max\{N^{2/3},N/q\})$ & open, top rung \\
\addlinespace[3pt]
\multicolumn{4}{@{}l}{\emph{class \cls{BP}, branching program, parallel
  (\cref{sec:BP})}}\\
\ref{BP0} & Thm  & repetition-free oblivious: $T\cdot(S+q\log_2 N)=\tilde\Omega(N)$ & proved \\
\ref{BP1} & Conj & $q=\sqrt N$, width $N^{O(1)}$: $T=\tilde\Omega(N)$ & open, weakest rung \\
\ref{BP2} & Conj & every $q\leq\sqrt N$: $T=\tilde\Omega(N)$ & open \\
\ref{BP3} & Conj & every $q$: $T=\tilde\Omega(\min\{N,N^{2}/q^{2}\})$ & open, top rung \\
\addlinespace[3pt]
\multicolumn{4}{@{}l}{\emph{general $k$ (in each model section)}}\\
\ref{RS2-k} & Conj & $S=O(1)$: $T=\omega(N^{(k-1)/k})$ & open \\
\ref{RS3-k} & Conj & $S=O(1)$: $T=N^{1-o(1)}$ & open \\
\ref{RS4-k} & Conj & every $S$: $T=\tilde\Omega(\max\{N^{(k-1)/k},NS^{-1/(k-2)}\})$ & open; upper bound attained \\
\ref{BS4-k} & Conj & every width: $T\cdot S^{1/(k-2)}=\tilde\Omega(N)$ & open \\
\ref{GS3-k} & Conj & every peak space: $\eCMCs=\tilde\Omega(N)$ & open \\
\ref{GP3-k} & Conj & every $q$: $\eCMCs=\tilde\Omega(\max\{N^{(k-1)/k},N/q\})$ & open \\
\ref{BP3-k} & Conj & every $q$: $T=\tilde\Omega(\min\{N,N^{k-1}/q^{k-1}\})$ & open \\
\bottomrule
\end{tabular}
\end{center}

\medskip\noindent
Classes \cls{US}, \cls{UP} and \cls{RP} carry no code: the first two are
settled by~\ref{F1} and \cref{p:onerou}, and the third reduces to
\cls{RS} by \cref{p:qcap,p:costfree}, with \cref{cor:depth} as its only
consequence.

% ---------------------------------------------------------------------
\section{A route through the problems}\label{sec:route}
% ---------------------------------------------------------------------

The ladders of \cref{p:ladders} order statements by logical strength, and
strength is not difficulty: \cref{BS1} is the weakest statement of its class
and stands behind the same obstruction as its strongest.  What governs
difficulty is whether any technique can see the class at all, and on that
axis the codes fall into seven stages.  Three principles generate the
order.  \emph{Restricting the algorithm buys more than weakening the
target}: all five proved statements restrict the algorithm and none weakens
a target.  \emph{Streaming models escape the barrier of \cref{r:barrier}
and the generic model does not}, which puts every oblivious statement ahead
of every adaptive one.  And \emph{a method that needs many outputs should be
given many outputs}, which puts \cref{sec:cfold} ahead of its own $C=1$
cases.

\begin{center}
\small
\begin{tabular}{@{}lp{0.34\textwidth}p{0.40\textwidth}@{}}
\toprule
stage & statements & what it takes \\
\midrule
0 & \ref{RS0}, \ref{RS5'}
  & counting inside the model; proved here \\
\addlinespace[2pt]
1 & \ref{BS0}, \ref{BP0}, \ref{GP0}, and the upper bound of \ref{RS4-k}
  & the streaming switching lemma \cite{Din20b} and the algorithm
    of~\cite{GNY25}; proved here, two caveats outstanding \\
\addlinespace[2pt]
2 & the oblivious curve at $k=3$; \ref{GP3} for oblivious adversaries at
    every peak space; \ref{BP3} for oblivious programs at fixed $q$
  & a cumulative-memory form of \cite{Din20b}; no new model and no new
    barrier \\
\addlinespace[2pt]
3 & \ref{BS4-C}, \ref{GS3-C}
  & one new estimate, then segmentation as in~\cite{Din20} and the
    transfer of \cref{r:transfer} \\
\addlinespace[2pt]
4 & \ref{RS1}--\ref{RS5}, \ref{RS2-k}, \ref{RS3-k}, \ref{RS4-k}
  & an argument in the functional graph of $f$ covering the non-lifting
    runs \\
\addlinespace[2pt]
5 & \ref{BS1}--\ref{BS4}, \ref{GS1}--\ref{GS4}, \ref{BS4-k}, \ref{GS3-k}
  & a way past the short-output barrier \\
\addlinespace[2pt]
6 & \ref{GP1}--\ref{GP3}, \ref{BP1}--\ref{BP3}, \ref{GP3-k}, \ref{BP3-k}
  & the same, and a parallel segmentation \\
\bottomrule
\end{tabular}
\end{center}

\subsection{Stage by stage}

\paragraph{Stage 0: what counting gives.}
\Cref{RS0,RS5'} are both union bounds over the operations of a run.
\Cref{RS5'} is the more useful of the two, because it subtracts from
\cref{RS4,RS5} everything a counting argument can reach and leaves a
sharply posed residue: the runs in which no $2$-collision is in hand when
the third preimage arrives.  Every algorithm in the literature, including
the one Hao, Huang and Liu isolate as an obstruction, lies on the side
already settled.

\paragraph{Stage 1: harvest.}
\Cref{BS0,BP0,GP0} are one lemma applied three times, and the upper bound
of \cref{RS4-k} is a reading of an existing algorithm.  What remains here
is bookkeeping rather than mathematics, in two places.  The three theorems
assume repetition-free schedules, and the multi-pass lemma covers a
schedule that repeats a fixed sequence but not one that repeats
arbitrarily; that gap should either be closed or recorded as a hypothesis.
And the algorithm behind \cref{RS4-k} uses auxiliary hash functions that
class \cls{RS} does not provide, so the substitution described after
\cref{RS4-k} has to be carried out.

\paragraph{Stage 2: finish the oblivious classes.}
Three targets, in increasing order of effort.  First, close the factor
$\sqrt S$ between \cref{BS0} and the $\tilde O(N/\sqrt S)$ algorithm inside
its own class; this is the smallest open quantity in the note and the only
one whose resolution needs no new technique in either direction.  Second,
remove the peak-space restriction from \cref{GP0}, which asks for a
cumulative-memory version of the streaming switching lemma: the conversion
of~\cite{BK23} is stated for proofs following Borodin and Cook, whereas
\cite{Din20b} argues through communication complexity, so this is an
adaptation and not a citation.  Third, close the factor $q$ between
\cref{BP0} and \cref{BP3}, whose hardest single point is $q = \sqrt N$ for
the reason given in \cref{r:parbase}.  None of the three changes the model,
and each is self-contained.

\paragraph{Stage 3: the $C$-fold statements.}
\Cref{BS4-C} at $C = \tilde\Omega(S)$ is the first assertion about
adaptive memory-bounded algorithms that current technique can plausibly
reach.  What has to be proved is the input property that segmentation
consumes: that a decision tree of height $h$ produces $k$ pairwise disjoint
$3$-collisions with probability exponentially small in $k$ as soon as
$h \ll (kN^{2})^{1/3}$, which is for triples what the corresponding
estimate for pairs is in~\cite{Din20} classically and in~\cite{HM21}
quantumly.  Given it, segmentation yields
\cref{BS4-C} and \cref{r:transfer} yields \cref{GS3-C} at no extra cost.
This is the stage at which a first theorem about class \cls{BS} should be
expected, and also the stage that would test whether the framing of this
note survives contact with a proof.

\paragraph{Stage 4: the register ladder.}
Begin at \cref{RS1}.  A success probability of $1-o(1)$ permits arguing
about typical executions rather than tail events; the qualitative form asks
for no exponent; and \cref{RS5'} has already reduced the task to the
non-lifting runs, where the quantity to be controlled is the factor $\ell$
of \cref{r:leverage}.  Any argument here must be sensitive to adaptivity,
since at $k=2$ the non-adaptive statement is \cref{RS0} and the adaptive one
is false by~\ref{F3}; the same sensitivity is what forced Dinur, Keller and
Marmor to leave the presampling technique behind in the preprocessing
model~\cite{DKM25}.  Then \cref{RS2,RS3,RS4,RS5} in that order, and then
the general-$k$ forms, whose slack $N^{1/k}$ shrinks as $k$ grows, so that
$k=3$ remains the easiest instance rather than merely the first.

\paragraph{Stage 5: the generic sequential wall.}
Classes \cls{BS} and \cls{GS} are blocked in their entirety by
\cref{r:barrier}, and no amount of weakening within a class moves them,
which is why their nominally weakest rungs are no easier than their
strongest.  Two ways to spend effort here rather than against it.  Prove the
$C$-fold statements first and then ask precisely what $C=1$ needs that
$C = \tilde\Omega(S)$ does not; the answer is a progress measure that
survives down to a single output, which is the open problem of~\cite{HHL26}
stated in this note's terms.  Or follow their route and modify the problem
until a measure exists, then argue the modification inessential; a version
of their construction whose recording oracle records collision pairs of $f$
itself, rather than tuples of an independent oracle, would speak directly
to \cref{BS4}.

\paragraph{Stage 6: the generic parallel classes.}
Last, and not only by inheritance.  Class \cls{BP} is the one class in which
the naive bound is false rather than unproved, by \cref{p:parbp}, so a
statement here has to be correctly shaped before it can be attacked, which
is why \cref{r:notconj} declines to conjecture beyond polynomial width.
The oblivious shadow of this class is stage 2 and is open; until that is
settled there is little reason to attempt \cref{BP1}.

\subsection{Two disagreements with the ladders}

The order above departs from \cref{p:ladders} in exactly two places, and
both are worth stating plainly.  \Cref{BS1,GS1} are called the weakest rungs
of their classes and should nonetheless not be attacked first: their
weakness is logical, not technical, and \cref{sec:cfold} offers a genuinely
easier target in the same class.  And \cref{GS3} should not be attempted
after \cref{BS4} but at the same time, since by \cref{r:transfer} a proof of
the second along the expected lines carries the first.  Of the codes in
\cref{sec:index}, five are theorems; of the rest, the three items of stage 2
are reachable with the tools now in hand, stage 3 needs one new estimate,
and everything from stage 5 on needs an idea that does not yet exist in the
literature.  A third disagreement is with the literature rather than with
the ladders: by \cref{r:proxies}, three settled questions resemble this one
closely enough to be mistaken for it, so a result in the advice model, or
about the decision version, or about comparison-based programs, advances no
stage above.

% ---------------------------------------------------------------------
\appendix
\section{Models of computation}\label{app:models}
% ---------------------------------------------------------------------

The model boxes of \cref{sec:R,sec:B,sec:G} are prose, and prose leaves
choices implicit that a proof then has to make.  This appendix gives all of
them as code-based games in the style of~\cite{BR06}.  Three things come out
of it.  The eight classes of \cref{sec:labels} become restrictions on the
type of a single function, so that what distinguishes them can be read
rather than recited.  The principle of deferred decisions, on which every
estimate of what a short computation can find depends, becomes the
lazy-sampling line of the execution game.  And the output normalisation that
segmentation arguments need becomes a wrapper whose cost is read off its
code, which is what \cref{sec:route} asks for at stage 3.

\subsection{Syntax}\label{app:syntax}

\begin{definition}[branching adversary]\label{d:adv}
Let $N,T,S \in \mathbb{N}$ and $\Sigma = \{0,1\}^{S}$.  A \emph{branching
adversary} over $\NN$ of length $T$ and space $S$ is a tuple
$\advers = (\sigma_{0}, \advers.\qry, \advers.\nxt, \advers.\outp)$ with
\[
  \begin{array}{ll}
  \sigma_{0} \in \Sigma
    & \text{the initial state},\\
  \advers.\qry \colon [T] \times \Sigma \to \NN
    & \text{the point queried},\\
  \advers.\nxt \colon [T] \times \Sigma \times \NN \to \Sigma
    & \text{the state update},\\
  \advers.\outp \colon [T] \times \Sigma \times \NN \to (\NN^{3})^{*}
    & \text{the statements emitted}.
  \end{array}
\]
All four are fixed before $f$ is drawn.
\end{definition}

Three features of the definition carry weight.  The functions are total and
deterministic, so an adversary is a finite object.  They receive the step
index, which is what makes the model levelled.  And $\advers.\outp$ emits
statements as a side effect of a transition rather than at the end, which is
what lets a proof speak of the statements emitted during an interval of
time, as \cref{sec:cfold} needs.

\begin{convention}[the step counter is free]\label{c:counter}
Passing $i$ to $\advers.\qry$ and $\advers.\nxt$ gives the adversary a
program counter that is not charged to $S$, which matches the levelled DAG,
where a node names its own level.  An adversary charged for its counter is
one whose state space must encode $i$, so the two conventions differ by an
additive $\lceil\log_{2}T\rceil$ in $S$.  Bounds stated with $S$ replaced by
$S + \log_{2}T$ hold under either, which is why that substitution appears
wherever a proof unions over $\Sigma$ at several time steps.
\end{convention}

\begin{convention}[coins]\label{c:coins}
A randomised branching adversary is a distribution over the objects of
\cref{d:adv}.  Since $f$ is drawn from a fixed distribution, for every
randomised adversary there is a deterministic one with the same resources
and at least the same success probability, obtained by fixing the coins to a
best value.  This is the remark that lets \cref{BS4} be stated for
deterministic programs without loss.
\end{convention}

\subsection{Execution}\label{app:exec}

\begin{figure}[t]
\begin{minipage}[t]{0.48\textwidth}
\begin{gbox}
\gname{$\mathsf{Game}~\Exec^{\advers}_{N,T}$}
\begin{pc}
\item $H \gets \emptyset$; \ $\sigma \gets \sigma_{0}$; \ $L \gets \varepsilon$
\item $\mathsf{cost} \gets 0$; \ $\mathsf{fresh} \gets 0$
\item \textbf{for} $i = 1$ \textbf{to} $T$ \textbf{do}
\item \quad $x \gets \advers.\qry(i,\sigma)$
\item \quad $y \gets \RO(x)$; \ $\mathsf{cost} \gets \mathsf{cost}+1$
\item \quad $L \gets L \,\|\, \advers.\outp(i,\sigma,y)$
\item \quad $\sigma \gets \advers.\nxt(i,\sigma,y)$
\item \textbf{return} $(L,H)$
\end{pc}
\end{gbox}
\end{minipage}\hfill
\begin{minipage}[t]{0.48\textwidth}
\begin{gbox}
\gname{$\mathsf{procedure}~\RO(x)$}
\begin{pc}
\item[\scriptsize 9:] \textbf{if} $H[x] = \bot$ \textbf{then}
\item[\scriptsize 10:] \quad $H[x] \getsr \NN$
\item[\scriptsize 11:] \quad $\mathsf{fresh} \gets \mathsf{fresh}+1$
\item[\scriptsize 12:] \textbf{return} $H[x]$
\end{pc}
\end{gbox}
\begin{gbox}
\gname{$\mathsf{procedure}~\Fin_{C}(L,H)$}
\begin{pc}
\item[\scriptsize 13:] \textbf{for} $x$ in $L$ with $H[x]=\bot$ \textbf{do}
\item[\scriptsize 14:] \quad $H[x] \getsr \NN$
\item[\scriptsize 15:] \textbf{return} $\mathsf{disj3}_{C}(L,H)$
\end{pc}
\end{gbox}
\end{minipage}
\caption{Execution of a branching adversary against a lazily sampled random
oracle, and the finalisation that decides the task of \cref{sec:cfold}.  At
$C=1$ the task is the one this note is about.}
\label{fig:exec}
\end{figure}

\Cref{fig:exec} is the whole model.  The table $H$ is the oracle, sampled
lazily; $L$ is the list of emitted statements; $\mathsf{cost}$ counts
charged steps and $\mathsf{fresh}$ counts points whose image has been
revealed.  The predicate $\mathsf{disj3}_{C}(L,H)$ returns $1$ iff $L$
contains $C$ entries $(x^{t}_{1},x^{t}_{2},x^{t}_{3})$, $t \in [C]$, whose
$3C$ points are pairwise distinct and which satisfy
$H[x^{t}_{1}] = H[x^{t}_{2}] = H[x^{t}_{3}]$ for every $t$.  The advantage
of $\advers$ is
\[
  \Adv^{\Col}_{N,C,T}(\advers) \;=\;
  \Pr\bigl[\,\Fin_{C}\bigl(\Exec^{\advers}_{N,T}\bigr) = 1\,\bigr],
\]
and the statements of this note read, for instance,
$\Adv^{\Col}_{N,1,T}(\advers) = O\bigl((T^{3}+1)N^{-2}\bigr)$ for
unbounded $\Sigma$, which is~\ref{F1}; and
$\Adv^{\Col}_{N,1,T}(\advers) \geq \tfrac12 \Rightarrow T\cdot S =
\tilde\Omega(N)$, which is \cref{BS4}.

\begin{remark}[what line 14 is for]\label{r:guess}
An adversary may emit points it never queried.  Sampling their images in
$\Fin$ rather than forbidding them keeps the model faithful and makes the
resulting term visible: a triple none of whose points was queried is a
$3$-collision with probability $N^{-2}$, which is the additive
$O(CN^{-2})$ that appears in~\ref{F1} and in the proof of \cref{BS0}.  An
argument that assumes emitted points are queried has to earn the
assumption, and \cref{l:ver} is how.
\end{remark}

\begin{proposition}[equivalence with the levelled DAG]\label{p:dag}
Branching adversaries of length $T$ and space $S$ are in bijection, up to
unreachable states, with the branching programs of \cref{sec:bp} of length
$T$ and width $2^{S}$ whose edges carry output statements.
\end{proposition}

\begin{proof}
Given $\advers$, take the node set $[T+1] \times \Sigma$, label $(i,\sigma)$
with $\advers.\qry(i,\sigma)$, and for each $y$ draw the edge
$(i,\sigma) \to (i+1,\advers.\nxt(i,\sigma,y))$ carrying
$\advers.\outp(i,\sigma,y)$; each level then has $2^{S}$ nodes.  Conversely,
index the nodes of each level of a program of width $W$ by
$\{0,1\}^{\lceil\log_{2}W\rceil}$ and read off the three functions.  The
computation paths correspond step by step, hence so do the emitted
statements and the success events.
\end{proof}

\begin{center}
\small
\begin{tabular}{@{}ll@{}}
\toprule
DAG & code \\
\midrule
level $i$ & the loop index, free by \cref{c:counter} \\
node at level $i$ & the pair $(i,\sigma)$ \\
node label & $\advers.\qry(i,\sigma)$ \\
out-edge labelled $y$ & $\advers.\nxt(i,\sigma,y)$ \\
output statement on an edge & $\advers.\outp(i,\sigma,y)$ \\
width $W$ & $|\Sigma| = 2^{S}$ \\
computation path on $f$ & the trace of $\mathsf{Game}~\Exec$ \\
\bottomrule
\end{tabular}
\end{center}

\subsection{The classes as restrictions of one type}\label{app:classes}

Everything below changes \cref{d:adv} in one place only, which is the point
of writing it this way.

\paragraph{Oblivious, repetition-free, non-adaptive.}
$\advers$ is \emph{oblivious} if $\advers.\qry(i,\sigma)$ is independent of
$\sigma$, so that $\advers.\qry$ factors through $[T]$, and
\emph{repetition-free} if that induced map is injective, equivalently if
$\mathsf{fresh} = T$ on every execution.  These are the hypotheses of
\cref{BS0,BP0,GP0}.  Obliviousness is also what \emph{non-adaptive} means
here: the schedule is fixed before $f$ is drawn and only the emitted
statements may depend on answers, which is the hypothesis of \cref{RS0}
transposed to class \cls{BS}.

\paragraph{Unbounded space.}
Taking $\Sigma = \{0,1\}^{*}$ and charging only $\mathsf{cost}$ gives class
\cls{US}, where~\ref{F1} is the whole story.

\paragraph{Parallel width $q$.}
Replace the two middle types by
\[
  \advers.\qry \colon [r] \times \Sigma \to \NN^{q},
  \qquad
  \advers.\nxt \colon [r] \times \Sigma \times \NN^{q} \to \Sigma ,
\]
run the loop for $r$ rounds calling $\RO$ on each coordinate, and set
$T = rq$.  The second type is the substantive one: a round branches on its
whole answer tuple, so an execution carries
$S + q\lceil\log_{2}N\rceil$ bits across the interior of a round though only
$S$ cross its boundary.  That is exactly the term charged in the proofs of
\cref{BP0,GP0}, and \cref{r:notconj} is the observation that it is charged
there and nowhere else that is proved.

\paragraph{Cumulative memory.}
Take $\Sigma = \{0,1\}^{*}$, insert into \cref{fig:exec} the accumulator
\[
  \textrm{5a:}\quad \cmcs \gets \cmcs + |\sigma| + \lceil\log_{2}N\rceil
\]
once per iteration, and charge $\cmcs$ in place of $(T,S)$; $\eCMCs$ is its
expectation over the coins of $\RO$ and $\CMCs$ its supremum.  The second
summand is the query charge discussed in \cref{sec:cmc}, and the accumulator
is what \cref{r:sustained} would replace by a count of the iterations at
which $|\sigma| \geq s$.

\paragraph{Register machines.}
Here the state is structured and the transition is restricted to five
operations, which makes the class an abstract model of computation in
Maurer's sense~\cite{Mau05}, the generic group model of Shoup~\cite{Sho97}
being the familiar instance.  Fix $w_{1},w_{2},\dots \in \NN$ before $f$ is
drawn.  The state is $\sigma = (Z,b,hd)$ with
$Z \in (\NN\cup\{\bot\})^{S}$, $b \in \{0,1\}^{\beta}$ and
$hd \in \mathbb{N}$, and the adversary is the single function
\[
  \advers.\stp \colon \{0,1\}^{\beta} \times (\NN \cup \{\bot\})
  \;\longrightarrow\; \mathsf{Op} \times \{0,1\}^{\beta},
\]
\[
  \mathsf{Op} \;=\;
  \{\Advance,\Next,\Copy,\Update,\Halt\} \times [S]^{3},
\]
which receives the value written by the previous operation and returns the
next operation with a new control state.  \Cref{fig:reg} is the execution.

\begin{figure}[t]
\begin{minipage}[t]{0.53\textwidth}
\begin{gbox}
\gname{$\mathsf{Game}~\Exec\text{-}\mathrm{REG}^{\advers}_{N,T}$}
\begin{pc}
\item $H \gets \emptyset$; \ $Z \gets (\bot,\dots,\bot)$; \ $b \gets b_{0}$
\item $hd \gets 1$; \ $v \gets \bot$; \ $\mathsf{cost} \gets 0$
\item \textbf{repeat}
\item \quad $(\mathsf{op},b) \gets \advers.\stp(b,v)$; \ $v \gets \bot$
\item \quad \textbf{if} $\mathsf{op} = \Advance(i,j)$ \textbf{then}
\item \quad\quad \textbf{if} $Z[i] = \bot$ \textbf{then abort}
\item \quad\quad $v \gets \RO(Z[i])$; \ $Z[j] \gets v$
\item \quad\quad $\mathsf{cost} \gets \mathsf{cost}+1$
\item \quad \textbf{if} $\mathsf{op} = \Next(j)$ \textbf{then}
\item \quad\quad $v \gets w_{hd}$; \ $Z[j] \gets v$; \ $hd \gets hd+1$
\item \quad\quad $\mathsf{cost} \gets \mathsf{cost}+1$
\item \quad \textbf{if} $\mathsf{op} = \Copy(i,j)$ \textbf{then} $Z[j] \gets Z[i]$
\item \quad \textbf{if} $\mathsf{op} = \Halt(i,j,k)$ \textbf{then}
\item \quad\quad \textbf{return} $\bigl((Z[i],Z[j],Z[k]),H\bigr)$
\item \textbf{until} $\mathsf{cost} = T$
\end{pc}
\end{gbox}
\end{minipage}\hfill
\begin{minipage}[t]{0.43\textwidth}
\small
The three features that give class \cls{RS} its content, none of which is
present in \cref{fig:exec}:

\smallskip
\emph{Contents.}  A register is written only at lines 7, 10 and 12, so it
holds a point of the closure of $\{w_{m}\}$ under $f$.  The oracle is never
called on a point of the adversary's own devising.  This is what
\cref{r:leverage} exploits.

\smallskip
\emph{Charging.}  $\mathsf{cost}$ counts $\Advance$ \emph{and} $\Next$, so
reaching the $m$-th sequence point costs $m$, while the oracle is called
only at line 7.  Cost and query count are different quantities.

\smallskip
\emph{Feedback.}  The adversary learns $v$ and remembers only what it puts
in $b$; the space is $S$ cells and $\beta$ control bits.
\end{minipage}
\caption{Class \cls{RS} as a restriction of \cref{fig:exec}: structured
state, five operations, cost charged for two of them.}
\label{fig:reg}
\end{figure}

\begin{remark}[the feedback channel is not optional]\label{r:feedback}
The prose model of \cref{sec:reg} once said only that the algorithm is a
function of what it was last told.  What it is told has to be part of the
definition, and the choice decides the class.  If $\advers.\stp$ receives
the value $v$, as in \cref{fig:reg}, then distinguished-point techniques are
available, since testing whether $v$ has a prescribed prefix is a
computation on $v$; parallel collision search~\cite{OW99} and the tradeoff
algorithms built on it are then legal, which is what makes $\tilde O(N/S)$
an upper bound for the class and hence what makes \cref{RS4} tight.  Were
the report instead the equality pattern of the register file, as in a
generic group model, distinguished points would be excluded and that upper
bound would fail.  The generous choice is also the right one for a lower
bound, and it is the one adopted throughout this note.
\end{remark}

\begin{center}
\small
\begin{tabular}{@{}llll@{}}
\toprule
class & state space & query type & charged \\
\midrule
\cls{US} & $\{0,1\}^{*}$ & one point & $\mathsf{cost}$ \\
\cls{BS} & $\{0,1\}^{S}$ & one point & $\mathsf{cost}$, $S$ \\
\cls{RS} & $(\NN\cup\{\bot\})^{S}\times\{0,1\}^{\beta}\times\mathbb{N}$
  & $Z[i]$, via $\Advance$ & $\mathsf{cost}$, $S$ \\
\cls{GS} & $\{0,1\}^{*}$ & one point & $\cmcs$ \\
\cls{UP} & $\{0,1\}^{*}$ & $q$-tuple & $\mathsf{cost}=rq$ \\
\cls{BP} & $\{0,1\}^{S}$ & $q$-tuple & $\mathsf{cost}=rq$, $S$ \\
\cls{RP} & as \cls{RS} & $q$ operations per round & $\mathsf{cost}$, $S$ \\
\cls{GP} & $\{0,1\}^{*}$ & $q$-tuple & $\cmcs$ \\
\bottomrule
\end{tabular}
\end{center}

\subsection{Two transforms}\label{app:transforms}

\begin{lemma}[deferred decisions]\label{l:lazy}
Let $\Exec^{\advers}_{N,T}[\mathrm{eager}]$ be \cref{fig:exec} with line 1
replaced by $f \getsr \mathrm{Fun}(\NN,\NN)$ and $\RO(x)$ by
\textbf{return} $f(x)$.  The two games induce the same distribution on
$(L,H)$, reading $H$ as the restriction of $f$ to the queried points.
Consequently, if $x_{i} \notin \{x_{1},\dots,x_{i-1}\}$ then $y_{i}$ is
uniform on $\NN$ and independent of $y_{1},\dots,y_{i-1}$, and the
subsequence of answers to fresh queries is i.i.d.\ uniform whatever the
adversary's strategy.
\end{lemma}

\begin{proof}
In either game the value returned for a point is drawn uniformly and once,
and $\advers$ is a function of the values returned so far, so the two
experiments couple step by step.  The consequence is line 10: at the moment
$y_{i}$ is returned for a fresh $x_{i}$ it has just been drawn uniformly,
and no earlier line inspected it.
\end{proof}

\Cref{l:lazy} is the whole probabilistic content of the estimates that
\cref{sec:route} calls for at stage 3, and it is worth saying what it does
not give.  Adaptivity is harmless for the distribution of answers and for
nothing else: the adversary chooses which points to query, so which
collisions exist among queried points is under its influence, and only the
count of fresh queries is protected.  This is also the reason
$\mathsf{fresh}$ and $\mathsf{cost}$ are separate counters in
\cref{fig:exec}: a tail estimate must be indexed by the former, and an
adversary that re-queries pays the latter for no new randomness, which is
the mechanism behind the multi-pass bound used to prove \cref{BS0}.

\begin{figure}[t]
\begin{minipage}[t]{0.53\textwidth}
\begin{gbox}
\gname{$\mathsf{adversary}~\Ver[\advers]$, at step $4(i-1)+1+m$}
\begin{pc}
\item \textbf{state} $\sigma' = (\sigma,\tau)$,
  $\tau \in \NN^{3}\cup\{\bot\}$
\item \textbf{if} $m = 0$ \quad \emph{// the original step}
\item \quad query $\advers.\qry(i,\sigma)$
\item \quad on answer $y$: \ $\tau \gets \advers.\outp(i,\sigma,y)$;
  \ $\sigma \gets \advers.\nxt(i,\sigma,y)$; emit nothing
\item \textbf{if} $m \in \{1,2,3\}$ \quad \emph{// verification}
\item \quad query $\tau[m]$, or a fixed point if $\tau = \bot$
\item \quad on answer $y$: emit $\tau$ if $m=3$ and $\tau \neq \bot$
\end{pc}
\end{gbox}
\end{minipage}\hfill
\begin{minipage}[t]{0.43\textwidth}
\small
The wrapper delays each statement by three steps and queries its three
points first.  It holds one pending triple, so the state grows by
$3\lceil\log_{2}N\rceil$ bits; when the wrapped adversary is a register
machine the triple is already in registers and the growth is nil.
\end{minipage}
\caption{The verification wrapper of \cref{l:ver}.  Its purpose is the
invariant, not efficiency.}
\label{fig:ver}
\end{figure}

\begin{lemma}[verification wrapper]\label{l:ver}
Let $\advers$ have resources $(T,S)$.  Then $\Ver[\advers]$ of
\cref{fig:ver} has length $4T$ and space $S + 3\lceil\log_{2}N\rceil$,
satisfies $\Adv^{\Col}_{N,C,4T}(\Ver[\advers]) =
\Adv^{\Col}_{N,C,T}(\advers)$, and emits every statement within three steps
of querying its three points.
\end{lemma}

\begin{proof}
Step $4(i-1)+1$ of the wrapper reproduces step $i$ of $\advers$ and its
answers are distributed identically, so the trace of $\sigma$ and the
multiset of statements agree; the interleaved queries are answered by $\RO$
and their answers discarded, so they influence neither $\sigma$ nor the
statements.  The invariant is immediate from lines 6 and 7, and $\tau$ holds
one triple.
\end{proof}

Two consequences.  With \cref{l:ver} in hand, line 14 of $\Fin$ never fires,
so the guessing term of \cref{r:guess} disappears at the cost of a factor
$4$ in length; this is the normalisation a segmentation argument needs, and
it is where the $C$-fold statements of \cref{sec:cfold} would begin.  And
the space cost is additive, so a bound of the form
$\mathsf{cost} \cdot S = \tilde\Omega(\cdot)$ is unaffected, whereas a bound
carrying $S$ in an exponent has to be restated for $S + O(\log N)$.

\subsection{Four conventions the prose leaves implicit}\label{app:pitfalls}

\begin{enumerate}[leftmargin=*,label=(\roman*)]\itemsep3pt
\item \emph{The step counter} is free, by \cref{c:counter}, so ``$S$ bits of
  memory'' means $S$ besides the counter.  A proof that unions over
  $\Sigma$ at each of several time steps pays $\log_{2}T$ anyway and should
  say $S + \log_{2}T$, as \cref{RS5} and \cref{sec:route} do.
\item \emph{Repeats} carry no randomness, by \cref{l:lazy}, so an estimate
  must be indexed by $\mathsf{fresh}$ and not by $\mathsf{cost}$.  The gap
  between the two is what separates adaptive cycle finding from the
  multi-pass streaming bound, and hence~\ref{F3} from \cref{BS0}.
\item \emph{Unqueried outputs} are permitted, and line 14 of $\Fin$ prices
  them rather than assuming them away.
\item \emph{Cost is not the query count} in class \cls{RS}, where $\Next$ is
  charged and calls no oracle.  A bound obtained from a query count becomes
  a bound on cost only after that gap is closed, which is why the proof of
  \cref{RS5'} counts charged operations rather than oracle calls.
\end{enumerate}

% ---------------------------------------------------------------------
\begin{thebibliography}{9}
\small

\bibitem{ABP18}
J.~Alwen, J.~Blocki, and K.~Pietrzak.
\newblock Sustained space complexity.
\newblock In \emph{EUROCRYPT 2018, Part II}, LNCS vol.~10821,
  pp.~99--130, Springer, 2018.

\bibitem{ACDW20}
Akshima, D.~Cash, A.~Drucker, and H.~Wee.
\newblock Time-space tradeoffs and short collisions in
  Merkle--Damg\aa{}rd hash functions.
\newblock In \emph{CRYPTO 2020, Part I}, LNCS vol.~12170, pp.~157--186,
  Springer, 2020.

\bibitem{Aks24}
Akshima.
\newblock Time-space tradeoffs for finding multi-collisions in
  Merkle--Damg\aa{}rd hash functions.
\newblock In \emph{ITC 2024}, LIPIcs vol.~304, pp.~9:1--9:22, Schloss
  Dagstuhl, 2024.

\bibitem{AS15}
J.~Alwen and V.~Serbinenko.
\newblock High parallel complexity graphs and memory-hard functions.
\newblock In \emph{STOC 2015}, pp.~595--603, ACM, 2015.

\bibitem{BC82}
A.~Borodin and S.~A. Cook.
\newblock A time-space tradeoff for sorting on a general sequential model
  of computation.
\newblock \emph{SIAM Journal on Computing} 11(2):287--297, 1982.

\bibitem{BFMUW87}
A.~Borodin, F.~Fich, F.~Meyer auf der Heide, E.~Upfal, and A.~Wigderson.
\newblock A time-space tradeoff for element distinctness.
\newblock \emph{SIAM Journal on Computing} 16(1):97--99, 1987.

\bibitem{BK16}
A.~Biryukov and D.~Khovratovich.
\newblock Equihash: asymmetric proof-of-work based on the generalized
  birthday problem.
\newblock In \emph{NDSS 2016}, Internet Society, 2016.  Journal version:
  \emph{Ledger} 2:1--30, 2017.

\bibitem{BK23}
P.~Beame and N.~Kornerup.
\newblock Cumulative memory lower bounds for randomized and quantum
  computation.
\newblock In \emph{ICALP 2023}, LIPIcs vol.~261, pp.~17:1--17:20; journal
  version \emph{ACM Transactions on Computation Theory} 17(3):20, 2025.

\bibitem{BR06}
M.~Bellare and P.~Rogaway.
\newblock The security of triple encryption and a framework for code-based
  game-playing proofs.
\newblock In \emph{EUROCRYPT 2006}, LNCS vol.~4004, pp.~409--426,
  Springer, 2006.

\bibitem{Brent80}
R.~P. Brent.
\newblock An improved Monte Carlo factorization algorithm.
\newblock \emph{BIT} 20:176--184, 1980.

\bibitem{BT24}
J.~Buzek and S.~Tessaro.
\newblock Collision resistance from multi-collision resistance for all
  constant parameters.
\newblock In \emph{CRYPTO 2024, Part V}, LNCS vol.~14924, pp.~429--458,
  Springer, 2024.

\bibitem{Din20}
I.~Dinur.
\newblock Tight time-space lower bounds for finding multiple collision
  pairs and their applications.
\newblock In \emph{EUROCRYPT 2020, Part I}, LNCS vol.~12105,
  pp.~405--434, Springer, 2020.

\bibitem{Din20b}
I.~Dinur.
\newblock On the streaming indistinguishability of a random permutation and
  a random function.
\newblock In \emph{EUROCRYPT 2020, Part II}, LNCS vol.~12106,
  pp.~433--460, Springer, 2020.

\bibitem{DKM25}
I.~Dinur, N.~Keller, and A.~Marmor.
\newblock Non-adaptive cryptanalytic time-space lower bounds via a
  Shearer-like inequality for permutations.
\newblock arXiv:2505.00894, 2025.

\bibitem{FO90}
P.~Flajolet and A.~M. Odlyzko.
\newblock Random mapping statistics.
\newblock In \emph{EUROCRYPT '89}, LNCS vol.~434, pp.~329--354, Springer,
  1990.

\bibitem{GNY25}
J.~Guo, W.~Nan, and Y.~Yao.
\newblock Revisiting time-space tradeoffs in collision search and decision
  problems.
\newblock In \emph{ASIACRYPT 2025}, LNCS vol.~16245, pp.~514--545,
  Springer, 2026.

\bibitem{HHL26}
Z.~Hao, Z.~Huang, and Q.~Liu.
\newblock On the need for (quantum) memory with short outputs.
\newblock arXiv:2602.23763, 2026.

\bibitem{HM21}
Y.~Hamoudi and F.~Magniez.
\newblock Quantum time-space tradeoff for finding multiple collision
  pairs.
\newblock In \emph{TQC 2021}, LIPIcs vol.~197, pp.~1:1--1:21, Schloss
  Dagstuhl, 2021.

\bibitem{JL09}
A.~Joux and S.~Lucks.
\newblock Improved generic algorithms for 3-collisions.
\newblock In \emph{ASIACRYPT 2009}, LNCS vol.~5912, pp.~347--363,
  Springer, 2009.

\bibitem{JT19}
J.~Jaeger and S.~Tessaro.
\newblock Tight time-memory trade-offs for symmetric encryption.
\newblock In \emph{EUROCRYPT 2019, Part I}, LNCS vol.~11476,
  pp.~467--497, Springer, 2019.

\bibitem{Mau05}
U.~Maurer.
\newblock Abstract models of computation in cryptography.
\newblock In \emph{Cryptography and Coding 2005}, LNCS vol.~3796,
  pp.~1--12, Springer, 2005.

\bibitem{OW99}
P.~C. van Oorschot and M.~J. Wiener.
\newblock Parallel collision search with cryptanalytic applications.
\newblock \emph{Journal of Cryptology} 12(1):1--28, 1999.

\bibitem{Pollard75}
J.~M. Pollard.
\newblock A Monte Carlo method for factorization.
\newblock \emph{BIT} 15(3):331--334, 1975.

\bibitem{RV22}
R.~D. Rothblum and P.~N. Vasudevan.
\newblock Collision-resistance from multi-collision-resistance.
\newblock In \emph{CRYPTO 2022, Part III}, LNCS vol.~13509,
  pp.~503--529, Springer, 2022.

\bibitem{Sho97}
V.~Shoup.
\newblock Lower bounds for discrete logarithms and related problems.
\newblock In \emph{EUROCRYPT '97}, LNCS vol.~1233, pp.~256--266,
  Springer, 1997.

\bibitem{STKT06}
K.~Suzuki, D.~Tonien, K.~Kurosawa, and K.~Toyota.
\newblock Birthday paradox for multi-collisions.
\newblock In \emph{ICISC 2006}, LNCS vol.~4296, pp.~29--40, Springer,
  2006.

\bibitem{Wag02}
D.~Wagner.
\newblock A generalized birthday problem.
\newblock In \emph{CRYPTO 2002}, LNCS vol.~2442, pp.~288--303, Springer,
  2002.

\bibitem{Yao94}
A.~C.-C. Yao.
\newblock Near-optimal time-space tradeoff for element distinctness.
\newblock \emph{SIAM Journal on Computing} 23(5):966--975, 1994.

\bibitem{Note2}
The author.
\newblock Many three-way collisions with bounded memory: the attackable
  case.
\newblock Companion note, this volume.

\end{thebibliography}

\end{document}
